\chapter{Integrales dobles y cálculo de áreas}

Presentamos algunos problemas relativos a las integrales dobles, cuyos principales elementos teóricos se presentaron en la sección \secref{integralesdobles}.
Se ilustrará la manera de plantear sumas de Riemann, el intercambio en el orden de integración y se parametrizará la región. Esto se hará básicamente, aplicando las fórmulas \ref{integralesdobles1} y \ref{integralesdobles2}.
\section{
Sumas de Riemmann
}
\begin{ejemplo}\normalfont
Evaluar mediante una suma de Riemann el volumen del sólido bajo la superficie $f(x,y)=6-\frac{1}{4}x^{3}-y^{2}$ sobre el plano $xy$ con $-1\leq x\leq 2,$ $0\leq y\leq 2.$

El rectángulo $R=\left[ -1,2\right] \times \left[ 0,2\right] $ se dividirá en $nm$ subrectángulos, $n$ para el intervalo $\left[ -1,2\right] $ y  $m$ para el intervalo $\left[0,2\right] .$ Asumiremos que los rectángulos en $x$ e $y,$ son iguales, respectivamente; esto es
$\Delta x=\frac{2-(-1) }{n}= \frac{3}{n},\,\,\Delta y=\frac{2-0}{m}= \frac{2}{m}.$
Por otra parte,


\begin{multline*}
x_{0}=-1,x_{1}=-1+\Delta x,x_{2}=-1+2\Delta x,\,\cdots, \\
x_{i}=-1+i\Delta x,\cdots,  x_{n}=-1+n\Delta x \\[0.1cm]
y_{0}=0,y_{1}=0+\Delta y,y_2=2\Delta y,\,\cdots ,y_{j}=j\Delta y,\,\cdots , \\ y_{m}=m\Delta y=m\left(\frac{2}{m}\right) = 2.
\end{multline*}



\begin{figure}[H]
\centering
\caption{$f(x,y)=6-\frac{1}{4}x^{3}-y^{2}.$}
\includegraphics[width=0.4\textwidth]{Img/T_69c.pdf} \hspace{5mm}
\includegraphics[width=0.4\textwidth]{Img/T_69b.pdf}

\includegraphics[width=0.4\textwidth]{Img/T_69a.pdf}
\end{figure}

En resumen $x_{i}=-1+\frac{3i}{n},$ $y_{j}=\frac{2j}{m}.$
Si $( x_{i}^{\ast },y_{j}^{\ast}) $ es un punto cualquiera en el rectángulo
$\left[ x_{i},x_{i+1}\right] \times \left[ y_{j},y_{j+1}\right],$
en particular consideraremos $\left( x_{i},y_{j}\right) ,$ entonces el volumen aproximado de paralelepípedo está dado por
\begin{eqnarray*}
f\left( x_{i},y_{j}\right) \Delta x\Delta y &=&\left( 6-\frac{1}{4}\left( -1+ \frac{3i}{n}\right) ^{3}-\left( \frac{2j}{m}\right) ^{2}\right) \frac{3}{n} \frac{2}{m} \\
&=& \frac{75}{2mn}-\frac{24j^{2}}{m^{3}n}+\frac{81i^{2}}{2mn^{3}}-\frac{81i^{3}}{2mn^{4}}-\frac{27i}{2mn^{2}}
\end{eqnarray*}
Sumando todos estos términos, obtendremos el volumen aproximado del sólido mencionado.
\begin{multline*}
\sum_{j=1}^m\sum_{i=1}^n f(x_i,y_j) \Delta x\Delta y
% = \sum_{j=1}^{m} \sum_{i=1}^{n}\left(\frac{75}{2mn}-\frac{24j^{2}}{m^{3}n}+\frac{81i^{2}}{2mn^{3}}-\frac{81i^{3}}{2mn^{4}}-\frac{27i}{2mn^{2}}\right)  \\
= \frac{209}{8}-\frac{4}{m^{2}}-\frac{27}{4n}-\frac{27}{8n^{2}}-\frac{12}{m}
\end{multline*}

La aproximación mejora si el número de rectángulos en que se divide $R$ es muy grande; de manera alternativa, si el ancho y el largo de dichos rectángulos tienden a cero; es por este motivo que se considera el siguiente límite:
\[
\lim_{n,m\to \infty }\left(\sum_{j=1}^{m}\sum_{i=1}^{n}f\left( x_{i},y_{j}\right) \Delta x\Delta y\right) = \frac{209}{8}.
\]
Evaluando la integral doble, se obtiene el mismo valor:
\[
\int_{-1}^{2}\int_{0}^{2}\left( 6-\frac{1}{4}x^{3}-y^{2}\right) \,dy\,dx
=\int_{-1}^{2}\left( \frac{28}{3}-\frac{1}{2}x^{3}\right) dx= \frac{209}{8}.
\]
Los gráficos, la suma de Riemann con su correspondiente límite y la integral doble se evalúan en \verb"Mathematica" con ayuda de las siguientes instrucciones:
\end{ejemplo}



\begin{tcolorbox}[breakable,enhanced,title=\bf Instrucción \verb"Mathematica", top=2pt,bottom=2pt,boxsep=0pt,before skip=2pt,after skip=0pt]
\begin{Verbatim}[formatcom=\letraverbatim]
f[x_,y_]:=6-x^3/4-y^2
Plot3D[f[x,y],{x,-1,2},{y,0,2}]
DiscretePlot3D[f[x,y],{x,-1,2,1/6},{y,0,2,1/8},
ExtentSize->Full,PlotStyle->Gray];
{xi,yj,dx,dy}={-1+3*i/n,2*j/m,3/n,2/m};
suma=Expand[Sum[f[xi,yj]*dx*dy,{i,1,n},{j,1,m}]]
Limit[suma,{n,m}->{Infinity,Infinity}]
Integrate[f[x,y],{x,-1,2},{y,0,2}]
\end{Verbatim}
\end{tcolorbox}



\begin{ejemplo}\normalfont
Usando sumas de Riemann, hallar el volumen del sólido bajo la superficie $f(x,y) =3-x^{2}y$ sobre el plano $xy,$ con $-3/2\leq x\leq 3/2,$
$0\leq y\leq 5/4.$

El rectángulo $\left[ -3/2,3/2\right] \times \left[ 0,5/4\right] $ se dividirá en $nm$ subrectángulos.
$n$ para el intervalo $\left[ -3/2,3/2\right] $ y $m$ para el intervalo $\left[ 0,5/4\right] $.




\begin{figure}[H]
\centering
\caption{$f(x,y)=3-x^{2}y$}
\includegraphics[width=0.4\textwidth]{Img/T_70a.pdf} \hspace{5mm}
\includegraphics[width=0.4\textwidth]{Img/T_70b.pdf}

\includegraphics[width=0.4\textwidth]{Img/T_70c.pdf}
\end{figure}

En este caso, la partición de los intervalos está dada por:
\[
\Delta x=\frac{\frac{3}{2}-\left(-\frac{3}{2}\right) }{n}= \frac{3}{n}\qquad \text{y}\qquad \Delta y=\frac{5/4-0}{m}= \frac{5}{4m}\text{.}
\]
Entonces, los puntos de partición son:
\[
x_{i}=-\frac{3}{2}+i\Delta x = -\frac{3}{2}+\frac{3i}{n}\qquad \text{y}\qquad y_{j}=j\Delta y = \frac{5j}{4m}\text{.}
\]
Evaluando la función $f(x, y)$ en estos puntos, se obtiene:
\[
f\left( x_{i},y_{j}\right) =3-x_{i}^{2}y_{j}=3-\left(\frac{3i}{n}-\frac{3}{2}\right)^{2}\left(\frac{5j}{4m}\right) = 3-\frac{45 j i^2}{4 m n^2}+\frac{45 j i}{4 m n}-\frac{45 j}{16 m}.
\]
Aplicando la suma de Riemann, el resultado es:
\begin{equation*}
	\sum_{i=1}^{n}\sum_{j=1}^{m}f\left( x_{i},y_{j}\right) \Delta x\Delta y = \frac{1215}{128}-\frac{225}{64 m n^2}-\frac{225}{128 m}-\frac{225}{64 n^2}.
\end{equation*}
Al tomar el límite de la suma de Riemann, se obtiene:
\[
\lim_{n,m\rightarrow \infty }\left( \sum_{i=1}^{n}\sum_{j=1}^{m}f(x_{i},y_{j}) \Delta x\Delta y\right) = \frac{1215}{128}
\]
Por otro lado, al evaluar la integral doble se obtiene el mismo valor:
\[
\int_{-\frac{3}{2}}^{\frac{3}{2}}\int_{0}^{\frac{5}{4}}\left( 3-x^{2}y\right) dydx=\int_{-\frac{3}{2}}^{\frac{3}{2}}\left[ 3y-\frac{1}{2}x^{2}y^{2}\right] _{y=0}^{y=\frac{5}{4}}dx=\frac{1215}{128}\text{.}
\]
Esta suma de Riemann, su correspondiente límite y la integral doble, se evalúan en \verb"Mathematica" con ayuda de las siguientes instrucciones:
\end{ejemplo}



\begin{tcolorbox}[breakable,enhanced,title=\bf Instrucción \verb"Mathematica", top=2pt,bottom=2pt,boxsep=0pt,before skip=2pt,after skip=0pt]
\begin{Verbatim}[formatcom=\letraverbatim]
f[x_,y_]:=3-x^2*y
Plot3D[f[x,y],{x,-3/2,3/2},{y,0,5/4}]
DiscretePlot3D[f[x,y],{x,-3/2,3/2,1/10,
{y,0,5/4,1/10},ExtentSize->Full,PlotStyle->Gray]
{xk,yj,dx,dy}={-3/2+3*k/n,5/4*j/m,3/n,5/(4*m)};
suma=Sum[f[xk,yj]*dx*dy,{k,1,n},{j,1,m}]
Limit[suma,{n,m}->{Infinity,Infinity}]
Integrate[f[x,y],{x,-3/2,3/2},{y,0,5/4}]
\end{Verbatim}
\end{tcolorbox}




\begin{ejemplo}\normalfont
Utilizando una suma de Riemann, evaluar la integral doble de $f(x,y)=x^2 \sin (2y)+y^2 \cos (x)$ sobre el rectángulo $[0,\pi/2]\times[0,\pi/2].$

\vskip1cm


\begin{figure}[H]
\centering
\caption{$f(x,y)=x^2 \sin (2y)+y^2 \cos (x)$}
\includegraphics[width=0.4\textwidth]{Img/T_68a.pdf} \hspace{5mm}
\includegraphics[width=0.4\textwidth]{Img/T_68b.pdf}

\includegraphics[width=0.4\textwidth]{Img/T_68c.pdf}
\end{figure}

Si los rectángulos en $x$ e $y$ son iguales, entonces
$\Delta x=\frac{\pi}{2n}$ y $\Delta y=\frac{\pi}{2m};$ por lo tanto,
$x_i=\frac{\pi i}{2n}$ y $y_j=\frac{\pi j}{2m}$. Por otra parte,
\[f(x_i,y_j)\Delta x\Delta y=\frac{\pi^4i^2\sin\left(\frac{\pi j}{m}\right)}{4mn^3}+\frac{\pi^4j^2\cos\left(\frac{\pi i}{2n}\right)}{4m^3n}\]
y
\begin{equation*}
	\resizebox{\linewidth}{!}{$
		\displaystyle
		\sum_{i=1}^n\sum_{j=1}^m f(x_i,y_j)\,\Delta x\,\Delta y
		=\frac{\pi^4(2m^2+3m+1)}{192m^2n}\left(\cot\frac{\pi}{4n}-1\right)
		+\frac{\pi^4(2n^2+3n+1)}{96mn^2}\cot\frac{\pi}{2m}.
		$}
\end{equation*}

entonces
\begin{equation*}
	\resizebox{\linewidth}{!}{$
		\displaystyle
		\int_0^{\frac{\pi}{2}}\!\int_0^{\frac{\pi}{2}}\big(x^2\sin(2y)+y^2\cos x\big)\,dy\,dx
		= \lim_{n,m\to\infty}\Big(\sum_{i=1}^n\sum_{j=1}^m f(x_i,y_j)\,\Delta x\,\Delta y\Big)
		= \frac{\pi^3}{12}.
		$}
\end{equation*}

Las figuras y cálculos anteriores se hicieron en \verb"Mathematica" ejecutando las siguientes instrucciones.
\end{ejemplo}



\begin{tcolorbox}[breakable,enhanced,title=\bf Instrucción \verb"Mathematica", top=2pt,bottom=2pt,boxsep=0pt,before skip=2pt,after skip=0pt]
\begin{Verbatim}[formatcom=\letraverbatim]
f[x_,y_]:=y^2*Cos[x]+x^2*Sin[2*y]
DiscretePlot3D[f[x,y],{x,0,Pi/2,Pi/35},
{y,0,Pi/2,Pi/35},ExtentSize->Full,PlotStyle->Gray]
{xi,yj,dx,dy}={Pi/2*i/n,Pi/2*j/m,Pi/(2*n),Pi/(2*m)};
Expand[f[xi,yj]*dx*dy]
suma=Simplify[Sum[f[xi,yj]*dx*dy,{i,1,n},{j,1,m}]]
Limit[suma,{n,m}->{Infinity,Infinity}]
Integrate[f[x,y],{y,0,Pi/2},{x,0,Pi/2}]
\end{Verbatim}
\end{tcolorbox}



Uno de los procedimientos que se utilizan con frecuencia para evaluar integrales dobles se ilustra a continuación.

\section{Cambio en el orden de integración}



\begin{enumerate}[leftmargin=0cm,itemindent=.75cm,labelwidth=\itemindent,labelsep=0cm,align=left,label={{\bf \arabic*.\,}}]
\item Evaluar la integral de la función $f(x,y)=xy$, sobre la región $D$ en el primer cuadrante del plano $xy$ acotada por las figuras de las expresiones $y=x,\,\,\,x^{2}+y^{2}=4$ y $x=2.$La intersección de la recta inclinada y la circunferencia corresponden en el primer cuadrante a $x=\sqrt{2}.$
\end{enumerate}


\begin{figure}[H]
\centering
\caption{$y=x,$ $x^{2}+y^{2}=4$}
\includegraphics[width=0.8\textwidth]{Fig/83.pdf}
\fuentepropia
\end{figure}

Para integrar en el orden $dxdy$, hay que tener en cuenta que hay una parte de la región donde $x$
recorre los puntos (de manera horizontal) desde la circunferencia hasta la recta $x=2$, con
$y \in [0,\sqrt{2}]$ y hay otra región de la región en la cual $x$ recorre los puntos desde la recta $y=x$
hasta la recta $x=2,$ con $y \in [\sqrt{2},2].$ 


En el orden $dydx$ se observa del gráfico que $y$ recorre los puntos (de manera vertical)
desde la circunferencia hasta la recta $y=x$, con $x\in[\sqrt{2},2].$



\begin{figure}[H]
\centering
\caption{$x\leq y \leq \sqrt{4-x^{2}}$}
\includegraphics[width=0.4\textwidth]{Fig/84.pdf} \hspace{5mm}
\includegraphics[width=0.4\textwidth]{Fig/85.pdf}
\fuentepropia
\end{figure}

Apoyándonos en el gráfico, la integral $\displaystyle\iint_{D}xydA$ se puede evaluar como sigue
\[\int_{0}^{\sqrt{2}}\int_{\sqrt{4-y^{2}}}^{2}xydxdy+\int_{\sqrt{2}}^{2}\int_{y}^{2}xydxdy=\int_{\sqrt{2}}^{2}\int_{\sqrt{4-x^{2}}}^{x}xydydx=1.\]



En coordenadas polares, la ecuación $x^2+y^2=4$ corresponde a $r=2$ y la recta $x=2$ se escribe como $r=2\sec \theta$.
Del gráfico se observa que $r$ recorre los puntos desde la circunferencia $r=2$ hasta la recta $r=2\sec \theta$.
La integral, usando la fórmula \ref{integraldoblepolares}, se evalúa de la siguiente manera:
\[\int_{0}^{\pi /4}\int_{2}^{2\sec \theta }r^{3}\sin \theta \cos\theta drd\theta =1.\]


\begin{figure}[H]
\centering
\caption{$2 \leq r \leq 2 \sec \theta.$}
\includegraphics[width=0.5\textwidth]{Fig/86.pdf}
\fuentepropia
\end{figure}

Estos cálculos se hicieron ejecutando las siguientes instrucciones:



\begin{tcolorbox}[breakable,enhanced,title=\bf Instrucción \verb"Mathematica", top=2pt,bottom=2pt,boxsep=0pt,before skip=2pt,after skip=0pt]
\begin{Verbatim}[formatcom=\letraverbatim]
Integrate[x*y,{x,Sqrt[2],2},{y,Sqrt[4-x^2],x}]
Integrate[x*y,{y,0,Sqrt[2]},{x,Sqrt[4-y^2],2}]+
Integrate[x*y,{y,Sqrt[2],2},{x,y,2}]
\end{Verbatim}
\end{tcolorbox}




El mismo valor se obtiene de la siguiente manera.



\begin{tcolorbox}[breakable,enhanced,title=\bf Parametrizando la región de integración, top=2pt,bottom=2pt,boxsep=0pt,before skip=2pt,after skip=0pt]
Un punto sobre la circunferencia tiene la forma {\small $P\left( 2\cos u, 2\sin u\right)$}, con {\small $u\in\left[0,\frac{\pi}{4}\right]$} y un punto sobre la recta $y=x$, fijando la primera coordenada, es de forma {\small $Q\left( 2\cos u, 2\cos u\right)$}. Por lo tanto, una combinación convexa de estos puntos parametriza la región de integración, es decir,

{\footnotesize \[
\mathbf{r}(u,v) = (1-v)P + vQ = \left\langle 2\cos u,2(1-v)\sin u+2v\cos u,0\right\rangle.\]}
Se agregó un $0$ en el tercer componente con el fin de evaluar el producto cruz.
Derivando parcialmente se obtienen los siguientes vectores:

{\footnotesize
\begin{align*}
\mathbf{r}_{u} &=  \left\langle -2\sin u, 2(1-v)\cos u-2v\sin u, 0\right\rangle \\
\mathbf{r}_{v} &= \left\langle 0, 2\cos u-2\sin u, 0\right\rangle \\
\mathbf{r}_{u}\times \mathbf{r}_{v} &=\left\langle 0, 0, -4(\cos u-\sin u) \sin u\right\rangle.
\end{align*}
}
Además, {\small \[xy = 4\left((1-v)\sin u+v\cos u\right)\cos u\]} y
{\small $\left\Vert \mathbf{r}_{u}\times \mathbf{r}_{v}\right\Vert =\left\vert 4( \sin u- \cos u) \sin u\right\vert,$}
de modo que la integral $\iint_D xy dA$ se convierte en
\\[0.1em]
\begin{center}
\small 
$\int_{0}^{1}\int_{0}^{\pi /4}16\left|\cos u((1-v)\sin u+v\cos u)(\sin u-\cos u)\sin u\right| dv du= 1.$
\end{center}
\end{tcolorbox}



\begin{enumerate}[start=2,leftmargin=0cm,itemindent=.75cm,labelwidth=\itemindent,labelsep=0cm,align=left,label={{\bf \arabic*.\,}}]
\item Invertir el orden de integración y evaluar la siguiente integral:
{\footnotesize \[\displaystyle\int_{0}^{1}\int_{\sqrt{y}}^{2-y}xy\, dxdy.\]}

La región de integración corresponde a la figura limitada en el primer cuadrante, por las figuras de las expresiones
$y=0$, $y=x^2$ y $y=2-x$.
\end{enumerate}

\begin{figure}[H]
	\centering
	\caption{$\sqrt{y}\leq x \leq 2-y $}
	
	\includegraphics[width=0.28\linewidth]{Fig/80.pdf}
	\includegraphics[width=0.28\linewidth]{Fig/81.pdf}
	\includegraphics[width=0.25\linewidth]{Fig/82.pdf}
	
	\fuentepropia
\end{figure}
%

Se puede observar que en el orden $dydx$, $y$ recorre (verticalmente) los puntos desde el eje $x$ hasta la parábola, en una parte, y hasta la recta inclinada en otra parte. Para el orden $dxdy$, $x$ recorre (horizontalmente) los puntos desde la parábola hasta la recta. Esto da origen a las siguientes integrales dobles en el orden que se requiere.

{\footnotesize
\[ \int_{0}^{1}\int_{\sqrt{y}}^{2-y}xydxdy =
\int_{0}^{1}\int_{0}^{x^{2}}xydydx+\int_{1}^{2}\int_{0}^{2-x}xydydx= \frac{7}{24}.\]
}
Estas integrales se evalúan con la ayuda de las siguientes instrucciones:



\begin{tcolorbox}[breakable,enhanced,title=\bf Instrucción \verb"Mathematica", top=2pt,bottom=2pt,boxsep=0pt,before skip=2pt,after skip=0pt]
\begin{Verbatim}[formatcom=\letraverbatim]
Integrate[x*y,{y,0,1},{x,Sqrt[y],2-y}]
Integrate[x*y,{x,0,1},{y,0,x^2}]+Integrate[x*y,
{x,1,2},{y,0,2-x}]
\end{Verbatim}
\end{tcolorbox}



Otra manera de evaluar la misma integral es la siguiente:



\begin{tcolorbox}[breakable,enhanced,title=\bf Parametrizando la región de integración, top=2pt,bottom=2pt,boxsep=0pt,before skip=2pt,after skip=0pt]
	
	La porción de parábola se parametriza por $\la \sqrt{t},t\ra,$ mientras que la recta se representa por $\la 2-t,t\ra ,$ para $t\in \lbrack 0,1].$
	Una combinación convexa de estas parametriza la región de interés, esto es
	{\small
	\[{\bf r}(t,u)=(1-u)\la 2-t,t\ra +u\la \sqrt{t},t\ra = \la (u-1)(t-2)+u\sqrt{t},t\ra\]}
	con $u\in \lbrack 0,1].$ De esta función se obtienen los siguientes vectores:
	{\small
	\[{\bf r}_{u} =\la t+\sqrt{t}-2,0\ra,{\bf r}_{t} =\la u+\frac{u}{2\sqrt{t}}-1,1\ra.\]}
	
	Los extendemos a $\mathbb{R}^{3},$ para obtener
	${\bf r}_{u}\times {\bf r}_{t}= \la 0,0,t+\sqrt{t}-2\ra.$
	De lo cual se sigue que $\left\Vert{\bf r}_{u}\times{\bf r}_{t}\right\Vert=\left\vert t+\sqrt{t}-2\right\vert.$
	Con la definición de la función se observa que $xy=(\left(u-1\right) \left( t-2\right) +\sqrt{t}u) t,$ entonces:
	\\[0.6em]
	\resizebox{\linewidth}{!}{$
		\iint_{R}xydA=\int_{0}^{1}\int_{0}^{1}\left((u-1)(t-2)+u\sqrt{t}\right)t\left|t+\sqrt{t}-2\right|\,du\,dt=\frac{7}{24}
		$}
\end{tcolorbox}



\begin{enumerate}[start=3,leftmargin=0cm,itemindent=.75cm,labelwidth=\itemindent,labelsep=0cm,align=left,label={{\bf \arabic*.\,}}]
	\item Determinar el centro de masa de una lámina delgada cuya densidad de área está dada por
$\delta (x,y)=x$ $kg/m^2$ y posee la forma de la región en el primer cuadrante interior a la circunferencia con ecuación $x^{2}+y^{2}=4$ y exterior a la circunferencia $x^{2}+y^{2}=2x.$
\end{enumerate}



Para plantear las integrales dobles utilizando coordenadas rectangulares, se debe tener en cuenta que en el orden $dxdy$, el área considerada se divide en tres porciones. En una de ellas $x$ toma valores desde el eje $y$ hasta la circunferencia de radio $1$.


\begin{figure}[H]
\centering
\caption{$\sqrt{x^2-2x}\leq y \leq \sqrt{4-x^2}$}
\includegraphics[width=0.8\textwidth]{Fig/87.pdf}
\fuentepropia
\end{figure}

En una segunda región $x$ toma valores desde la circunferencia de radio $1$ hasta la circunferencia mayor y en una tercera región $x$ recorre los puntos desde el eje $y$ hasta la circunferencia de radio $2$. En el orden $dydx$, $y$ recorre los puntos desde la circunferencia de radio $1$ hasta la circunferencia de radio $2$ con $x\in [0,2].$

En coordenadas polares la ecuación de la circunferencia $x^2+y^2=4$ se reescribe como $r=2$
y la ecuación $x^2+y^2=2x$ corresponde a $r^2=2r\cos \theta,$ es decir, $r=2 \cos \theta.$ Es por esto que $r$
recorre los puntos desde la circunferencia $r=2 \cos \theta$ hasta $r=2;$
es decir, la región de interés se describe con $2\cos \theta \leq r \leq 2$ y $0\leq  \theta \leq \pi/2$.



\begin{figure}[H]
\centering
\caption{$2\cos \theta \leq r \leq 2$}
\includegraphics[width=0.4\textwidth]{Fig/88.pdf} \hspace{5mm}
\includegraphics[width=0.4\textwidth]{Fig/89.pdf}

\includegraphics[width=0.4\textwidth]{Fig/90.pdf}
\fuentepropia
\end{figure}

Estas figuras ayudan a establecer que la masa $M$ de la lámina en el orden $dxdy$ está dada por
\[
M=\int_{0}^{1}\int_{0}^{1-\sqrt{1-y^{2}}}xdxdy+\int_{0}^{1}\int_{1+\sqrt{1-y^{2}}}^{\sqrt{4-y^{2}}}xdxdy+\int_{1}^{2}\int_{0}^{\sqrt{4-y^{2}}}xdxdy
\]

En el orden $dydx$, la integral correspondiente es
\[
M=\int_{0}^{2}\int_{\sqrt{2x-x^{2}}}^{\sqrt{4-x^{2}}}xdydx=\frac{8}{3}-\frac{\pi}{2}.\]
y en coordenadas polares la masa está dada por
\[M=\ds\int_{0}^{\pi/2}\int_{2\cos\theta}^2\left(r\cos\theta\right)rdrd\theta=\frac{8}{3}-\frac{\pi}{2}.\]
Mientras que los momentos con respecto al eje  $y$ y al eje $x$ están dados por
\begin{eqnarray*}
\mu {y} &=& \int_{0}^{2}\int_{\sqrt{2x-x^{2}}}^{\sqrt{4-x^{2}}}x^2dydx=\int_{0}^{\pi /2}\int_{2\cos \theta }^{2}r^{3}\cos ^{2}\theta drd\theta = \frac{3}{8}\pi \\
\mu {x} &=& \int_{0}^{2}\int_{\sqrt{2x-x^{2}}}^{\sqrt{4-x^{2}}}xydydx=\int_{0}^{\pi /2}\int_{2\cos \theta }^{2}r^{3}\cos \theta \sin \theta drd\theta = \frac{4}{3}.
\end{eqnarray*}
Por lo tanto, el centro de masa de la lámina es
\[\left( \overline{x},\overline{y}\right) = \left( \frac{\mu _{y}}{M},\frac{\mu _{x}}{M}\right) = \left( \frac{\frac{3}{8}\pi }{\frac{8}{3}-\frac{1}{2}\pi },\frac{\frac{4}{3}}{\frac{8}{3}-\frac{1}{2}\pi }\right) = \left( \frac{9\pi }{64-12\pi },\frac{8}{16-3\pi}\right). \]

Los cálculos para obtener los momentos y la masa se hacen usando las siguientes instrucciones:



\begin{tcolorbox}[breakable,enhanced,title=\bf Instrucción en \verb"Mathematica", top=2pt,bottom=2pt,boxsep=0pt,before skip=2pt,after skip=0pt]
\begin{Verbatim}[formatcom=\letraverbatim]
Integrate[x,{x,0,2},{y,Sqrt[2*x-x^2],Sqrt[4-x^2]}]
Integrate[x*y,{x,0,2},{y,Sqrt[2*x-x^2],Sqrt[4-x^2]}]
Integrate[x^2,{x,0,2},{y,Sqrt[2*x-x^2],Sqrt[4-x^2]}]
Integrate[r^2*Cos[t],{t,0,Pi/2},{r,2*Cos[t],2}]
\end{Verbatim}
\end{tcolorbox}



%\clearpage

\begin{enumerate}[start=4,leftmargin=0cm,itemindent=.75cm,labelwidth=\itemindent,labelsep=0cm,align=left,label={{\bf \arabic*.\,}}]
	
\item
Determinar las coordenadas $(\overline{x},\overline{y})$ del centro de masa de una lámina determinada por la región en el primer cuadrante acotada por las figuras de las ecuaciones $x^2+y^2=4$ y $x^2+y^2=2y,$
cuya densidad de masa en cualquier punto $(x,y)$ es $\delta (x,y)=xy^2$. La forma de la lámina se presenta en la figura \ref{fig510}.
\end{enumerate}


\begin{figure}[H]
\centering
\caption{$2\sin \theta \leq r \leq 2$}
\label{fig510}
\includegraphics[width=0.8\textwidth]{Fig/91.pdf}
\fuentepropia
\end{figure}

Con el propósito de plantear las integrales dobles que determinan la masa y los demás elementos requeridos, observemos en  los siguientes gráficos la manera como se recorre la región de integración.


En el orden $dxdy$, $x$ recorre los puntos desde la circunferencia menor hasta la mayor, para $y \in [0,2]$. En el orden $dydx$, la región se divide en tres porciones. En una parte $y$ recorre los puntos desde el eje $x$ hasta la circunferencia menor con $y \in [0,1]$; en una segunda parte $y$ recorre los puntos desde la circunferencia menor hasta la circunferencia mayor, con $y \in [1,2]$, y en la tercera parte $y$ recorre los puntos desde el eje $x$ hasta la circunferencia mayor con $y \in [0,\sqrt{3}]$. También se debe observar que de la ecuación $x^2+y^2=2y$ se desprende $y=1 \pm \sqrt{1-x^2}.$


\begin{figure}[H]
\centering
\caption{$2\sin \theta \leq r \leq 2$}

\includegraphics[width=0.30\linewidth]{Fig/92.pdf}
\includegraphics[width=0.26\linewidth]{Fig/93.pdf}
\includegraphics[width=0.26\linewidth]{Fig/94.pdf}
\fuentepropia
\end{figure}

En coordenadas polares se satisface que $2\sin \theta \leq r \leq 2$ para $0 \leq \theta \leq \pi/2$. Por tanto, la masa $M$ de la lámina se puede hallar evaluando cualquiera de las siguientes integrales.
%
\[
\resizebox{\textwidth}{!}{$
\begin{aligned}
M &= \int_{0}^{1} \int_{0}^{1-\sqrt{1-x^{2}}}xy^{2} \,dy\,dx 
+ \int_{0}^{1} \int_{1+\sqrt{1-x^{2}}}^{\sqrt{4-x^{2}}} xy^{2} \,dy\,dx 
+ \int_{1}^{2} \int_{0}^{\sqrt{4-x^{2}}} xy^{2} \,dy\,dx \\
&= \int_{0}^{2} \int_{\sqrt{2y - y^{2}}}^{\sqrt{4 - y^{2}}} xy^{2} \,dx\,dy 
= \int_{0}^{\frac{\pi}{2}} \int_{2\sin\theta}^{2} r^{4} \sin^{2}\theta \cos\theta \,dr\,d\theta 
= \frac{4}{3}
\end{aligned}
$}
\]

%
Los momentos con respecto al eje $x$, $\mu_{x}$ y, con respecto al eje $y$, $\mu_{y}$,
se pueden determinar de la siguiente manera:

$\mu_{x}=\int_{0}^{2}\int_{\sqrt{2y-y^{2}}}^{\sqrt{4-y^{2}}}xy^{3}dxdy=\int_{0}^{\frac{\pi }{2}}\int_{2\sin \theta }^2r^{5}\cos \theta \sin^{3}\theta drd\theta=\frac{8}{5}$

$\mu_{y}= \int_{0}^{2}\int_{\sqrt{2y-y^{2}}}^{\sqrt{4-y^{2}}}x^{2}y^{2}dxdy=\int_{0}^{\pi /2}\int_{2\sin \theta }^{2} r^5 \cos^2 \theta\sin^2 \theta drd\theta =\frac{25}{48}\pi .$


Con lo anterior, se concluye que el centro de masa está dado por
\[(\overline{x},\overline{y})= \left( \frac{\frac{25}{48}\pi }{\frac{4}{3}},\frac{\frac{8}{5}}{\frac{4}{3}}\right) = \left( \frac{25\pi }{64},\frac{6}{5}\right).\]


El cálculo de la masa se puede evaluar mediante el uso de las siguientes indicaciones.



\begin{tcolorbox}[breakable,enhanced,title=\bf Instrucción en \verb"Mathematica", top=2pt,bottom=2pt,boxsep=0pt,before skip=2pt,after skip=0pt]
\begin{Verbatim}[formatcom=\letraverbatim]
Integrate[x*y^2,{y,0,2},{x,Sqrt[2*y-y^2],
Sqrt[4-y^2]}]
Integrate[r^4*Sin[t]^2*Cos[t],{t,0,Pi/2},
{r,2*Sin[t],2}]
\end{Verbatim}
\end{tcolorbox}



\begin{enumerate}[start=5,leftmargin=0cm,itemindent=.75cm,labelwidth=\itemindent,labelsep=0cm,align=left,label={{\bf \arabic*.\,}}]
\item
Si $f$ es continua en $R$, plantear la integral $\ds\iint_R f(x,y)dA$, siendo $R$ la región en el primer cuadrante acotada por las figuras de las ecuaciones $x^2+y^2=4,$ $y= \sqrt{3}\,x$ y $y=1.$
En coordenadas polares $y=1$ se transforma en $r= \csc \theta$, $y= \sqrt{3}\,x$ equivale a $\theta = \pi /3 $ y la ecuación $x=\sqrt{4-y^2}$ corresponde a un arco de circunferencia de radio $2$ con centro en el origen.
\end{enumerate}


\begin{figure}[H]
\centering
\caption{$\csc \theta \leq r \leq 2,$ $\frac{\pi}{6} \leq\theta \leq \frac{\pi}{3}$}
\includegraphics[width=0.5\textwidth]{Fig/95.pdf}
\fuentepropia
\end{figure}

En los siguientes gráficos se muestran las diversas maneras de recorrer la región $R$. Para integrar en el orden $dxdy$ se debe tener en cuenta que $x$ recorre los puntos desde la recta $y=\sqrt{3}\, x$ hasta la circunferencia. Para el orden $dydx$, en una parte $y$ recorre puntos desde la recta $y=1$ hasta la recta $y=\sqrt{3}\, x,$ y en otra parte $y$ recorre los puntos desde $y=1$ hasta la circunferencia.

% REDUCCIÓN DEL ESPACIO VERTICAL SUPERIOR
\vspace*{1.0em} 
\begin{center} 
	% Usamos un parbox principal para contener y alinear todo el bloque
	\parbox{\linewidth}{
		\centering
		
		% LEYENDA (encima de la figura)
		\footnotesize
		\textbf{Figura 13.} $\csc \theta \leq r \leq 2,\quad \frac{\pi}{6} \leq \theta \leq \frac{\pi}{3}.$
		 % Espacio mínimo entre leyenda e imágenes
		
		% LAS TRES IMÁGENES
		\includegraphics[width=0.28\linewidth]{Fig/96.pdf}
		\includegraphics[width=0.28\linewidth]{Fig/97.pdf}
		\includegraphics[width=0.28\linewidth]{Fig/98.pdf}
		
		% NOTA DE ELABORACIÓN
		\fuentepropia
	}
\end{center}
% REDUCCIÓN DEL ESPACIO VERTICAL INFERIOR (antes del siguiente párrafo)
\vspace*{1.2em} 

En coordenadas polares, $r$ toma valores desde la recta horizontal hasta la circunferencia con $\theta \in [\frac{\pi}{6},\frac{\pi}{3}].$
De acuerdo con lo anterior, se pueden plantear las integrales, como se muestra a continuación:


% REDUCCIÓN DEL ESPACIO DE LA ECUACIÓN
%\setlength{\abovedisplayskip}{3pt} 
%\setlength{\belowdisplayskip}{3pt}
\[
\large
\resizebox{\textwidth}{!}{$
	\begin{array}{rcl}
		\ds\iint_R f(x,y)\,dA &=& \int_{\frac{1}{\sqrt{3}}}^{1}\int_{1}^{\sqrt{3}x}f(x,y)\,dy\,dx 
		+ \int_{1}^{\sqrt{3}}\int_{y}^{\sqrt{4-x^{2}}}f(x,y)\,dy\,dx \\[10pt]
		&=& \int_{1}^{\sqrt{3}}\int_{y/\sqrt{3}}^{\sqrt{4-x^{2}}}f(x,y)\,dx\,dy 
		= \int_{\pi/6}^{\pi/3}\int_{\csc\theta}^{2}f(r,\theta)\,r\,dr\,d\theta
	\end{array}
	$}
\]

%\vspace*{1.2em} 
\begin{enumerate}[start=6,leftmargin=0cm,itemindent=.75cm,labelwidth=\itemindent,labelsep=0cm,align=left,label={{\bf \arabic*.\,}}]
	\item
Escribe en coordenadas rectangulares la siguiente integral: \[\ds\int_{0}^{\pi/4}\int_{0}^{2\sqrt{2}}r^2drd\theta.\]
\end{enumerate}


Los límites de integración son $r=0,$ $r=2\sqrt{2},$ $\theta=0,$ $\theta= \pi/4,$
esto determina una región en el primer cuadrante limitada por las rectas $y=0,$
$y=x,$ y la circunferencia con ecuación \[x^2+y^2=\left(2\sqrt{2}\right)^2=8.\]
Se usó un cambio a coordenadas polares $x= r \cos \theta$  y $y= r \cos \theta$.


\begin{figure}[H]
\centering
\caption{$0 \leq r \leq 2\sqrt{2}.$}
\includegraphics[width=0.8\textwidth]{Fig/99.pdf}
\fuentepropia
\end{figure}

Al evaluar la integral en el orden $dydx,$
la región de integración se divide en dos, en una parte $y$ varía desde el eje $x$
hasta la recta $y=x$ con $x \in[0,2],$ en la otra parte $y$ varía desde el eje $x$
hasta la circunferencia, con $x \in[2,2\sqrt{2}].$
Al integrar en  orden $dxdy$, el $x$ recorre (horizontalmente) los puntos desde la recta hasta la circunferencia, con $y\in[0,2]$.

\begin{figure}[h!]
	\centering
	\caption{$y \leq x \leq \sqrt{8 - x^2}.$}
	\includegraphics[width=0.3\linewidth]{Fig/100.pdf}
	\includegraphics[width=0.3\linewidth]{Fig/101.pdf}
	\includegraphics[width=0.3\linewidth]{Fig/102.pdf}
	\fuentepropia
\end{figure}

Por lo señalado antes, la integral en mención se reescribe como
\begin{multline*}
\displaystyle\int_{0}^{\pi /4}\int_{0}^{2\sqrt{2}}r^2drd\theta = \int_{0}^2\int_{0}^{x}\sqrt{x^2+y^2}dydx + \\ \int_{2}^{2\sqrt{2}}\int_{0}^{\sqrt{8-x^2}}\sqrt{x^2+y^2}dydx = \int_0^2 \int _y ^{\sqrt{8-y^2}} \sqrt{x^2+y^2}\, dxdy = \frac{4\sqrt{2} \pi}{3}
\end{multline*}

Estas integrales se evaluaron por medio de las siguientes instrucciones:

%

\begin{tcolorbox}[breakable,enhanced,title=\bf Instrucción en \verb"Mathematica", top=2pt,bottom=2pt,boxsep=0pt,before skip=2pt,after skip=0pt]
\begin{Verbatim}[formatcom=\letraverbatim]
Integrate[Sqrt[x^2+y^2],{y,0,2},{x,y,Sqrt[8-y^2]}]
Integrate[Sqrt[x^2+y^2],{x,0,2},{y,0,x}]+
Integrate[Sqrt[x^2+y^2],{x,2,2*Sqrt[2]},
{y,0,Sqrt[8-x^2]}]
Integrate[r^2,{\[Theta],0,Pi/4},{r,0,2*Sqrt[2]}]
\end{Verbatim}
\end{tcolorbox}




\begin{enumerate}[start=7,leftmargin=0cm,itemindent=.75cm,labelwidth=\itemindent,labelsep=0cm,align=left,label={{\bf \arabic*.\,}}]
\item 
Evaluar la integral $\ds \iint_{R} |y-x|dA,$ donde $R$ es la región acotada por las figuras de las ecuaciones $y=2-x^{2}$ y $y=x^{2}+2x-2.$


La función $|y-x|$ se expresa como

\[|y-x|= \left \{\begin{array}{cc} y-x & \text{si\,\,} y\geq x \\ x-y & \text{si\,\,} y<x\end{array} \right.\]
\end{enumerate}


\begin{figure}[H]
\centering
\caption{$x^2+2x-2 \leq y \leq 2-x^2$}
\includegraphics[width=0.8\textwidth]{Fig/103.pdf}
\fuentepropia
\end{figure}

El planteamiento de la integral doble en coordenadas rectangulares requiere hacer algunas consideraciones. De la expresión $y=2-x^2$ se obtiene $x= \pm \sqrt{2-y},$ mientras que de la ecuación  $y=x^2+2x-2$ se obtiene que $x=-1\pm \sqrt{y+3}$. Ahora se ilustra gráficamente la región de integración.

\begin{figure}[h!]
	\centering
	\caption{$x^2 + 2x - 2 \leq y \leq 2 - x^2$}
	\includegraphics[width=0.48\linewidth]{Fig/104.pdf}
	\includegraphics[width=0.38\linewidth]{Fig/105.pdf}
	\fuentepropia
\end{figure}

En el orden $dydx$ la región $R$ se divide en dos partes, en una de ellas $y$ recorre los puntos desde la recta $y=x$ hasta la curva $y=2-x^2,$ para $x\in [-2,1].$
En la otra región $y$ recorre los puntos desde la curva
$y=x^2+2x-1$ hasta la recta $y=x$, con $x\in [-2,1].$ En este caso, la integral es
\[\ds \int_{-2}^{1}\int_{x}^{2-x^{2}}\left( y-x\right)
dydx+\int_{-2}^{1}\int_{x^{2}+2x-2}^{x}\left( x-y\right) dyd = \frac{81}{10}\]

Con las siguientes instrucciones se evalúa esta integral.



\begin{tcolorbox}[breakable,enhanced,title=\bf Instrucción en \verb"Mathematica", top=2pt,bottom=2pt,boxsep=0pt,before skip=2pt,after skip=0pt]
\begin{Verbatim}[formatcom=\letraverbatim]
Integrate[y-x,{x,-2,1},{y,x,2-x^2}]+Integrate[x-y,
{x,-2,1},{y,x^2+2*x-2,x}]
\end{Verbatim}
\end{tcolorbox}



Para el orden $dxdy$, $R$ queda dividida en cuatro regiones de acuerdo con el recorrido horizontal que hace $x$.

\begin{itemize}[label=$\bullet$\,\,]
\item  $y\in[1,2]$ se satisface que $-\sqrt{2-y} \leq x \leq \sqrt{2-y}.$
\item Para $y\in[-2,1],$ hay dos regiones: en una parte, $-\sqrt{2-y} \leq x \leq y,$
mientras que en la otra región $y \leq x \leq 1+\sqrt{3+y}.$
\item Para $y\in[-3,-2]$ se cumple que $-1-\sqrt{1+y} \leq x \leq -1+\sqrt{1+y}.$
\end{itemize}
Por lo señalado antes, la integral en este caso es la siguiente:
\begin{multline*}
\int_{1}^{2}\int_{-\sqrt{2-y}}^{\sqrt{2-y}}(y-x)dxdy+\int_{-2}^{1}\int_{-\sqrt{2-y}}^{y}(y-x) dxdy \\
+\int_{-2}^{1}\int_{y}^{-1+\sqrt{3+y}}(x-y)dxdy
+\int_{-3}^{-2}\int_{-1-\sqrt{3+y}}^{-1+\sqrt{3+y}}(x-y)dxdy=\frac{81}{10}.
\end{multline*}
Las integrales anteriores se evalúan con las siguientes instrucciones:



\begin{tcolorbox}[breakable,enhanced,title=\bf Instrucción en \verb"Mathematica", top=2pt,bottom=2pt,boxsep=0pt,before skip=2pt,after skip=0pt]
\begin{Verbatim}[formatcom=\letraverbatim]
Integrate[y-x,{y,1,2},{x,-Sqrt[2-y],Sqrt[2-y]}]+
Integrate[y-x,{y,-2,1},{x,-Sqrt[2-y],y}]+Integrate[x-y,
{y,-2,1},{x,y,-1+Sqrt[3+y]}]+Integrate[x-y,{y,-3,-2},
{x,-1-Sqrt[3+y],-1+Sqrt[3+y]}]
\end{Verbatim}
\end{tcolorbox}



\begin{enumerate}[start=8,leftmargin=0cm,itemindent=.75cm,labelwidth=\itemindent,labelsep=0cm,align=left,label={{\bf \arabic*.\,}}]
	\item 
Evaluar $\ds \iint R x dA,$ donde $R$ es la región acotada por las figuras de las ecuaciones $x^2+y^2=4$
y $y=1-(1+\sqrt{3})x$ con  $x \geq 0.$ \\[0.1cm]

De la ecuación de la recta se sigue que
$x=\frac{1}{2}\left(1-\sqrt{3}\right)(y-1),$ un punto de corte de esta línea con la circunferencia es $(1,-\sqrt{3}).$
\end{enumerate}


\begin{figure}[H]
\centering
\caption{$x^2+y^2=4$}
\includegraphics[width=0.8\textwidth]{Fig/106.pdf}
\fuentepropia
\end{figure}

Para evaluar la integral en el orden $dydx$ se debe tener en cuenta dos regiones: en una de estas $y$ recorre los puntos desde la recta inclinada hasta la circunferencia, con $x \in [0,1],$ en la otra región $y$ recorre puntos desde la parte negativa de la circunferencia hasta la parte positiva de la misma, con $x \in [1,2].$
Al considerar el orden  $dxdy$ tendremos dos regiones, en una de ellas $x$ recorre los puntos desde el eje $y$ hasta la circunferencia,  con $y \in [1,2]$; en la otra región $x$ recorre los puntos desde la recta inclinada hasta la circunferencia, con $y \in [-\sqrt{3},1]$.


\begin{center}
	{\small\textbf{Figura:} $x^2+y^2\leq 4,\quad 1-(1+\sqrt{3})x \geq y$}
	
	
	
	\begin{minipage}{0.29\linewidth}
		\centering
		\includegraphics[width=\linewidth]{Fig/107.pdf}
	\end{minipage}
	\hfill
	\begin{minipage}{0.29\linewidth}
		\centering
		\includegraphics[width=\linewidth]{Fig/108.pdf}
	\end{minipage}
	\hfill
	\begin{minipage}{0.29\linewidth}
		\centering
		\includegraphics[width=\linewidth]{Fig/109.pdf}
	\end{minipage}
	
	
	
	\fuentepropia
\end{center}

Por lo señalado antes, la integral se puede evaluar como sigue.
{\small
\begin{align*}
	\ds \iint R x dA &= \int_{-\sqrt{3}}^{1}\int_{\frac{1}{2}\left(1-\sqrt{3}\right)(y-1)}^{\sqrt{4-y^{2}}}\!\!xdxdy + \int_1^2\int_{0}^{\sqrt{4-y^{2}}}\!\!xdxdy \\
	&= \int_{0}^{1}\int_{1-\left( 1+\sqrt{3}\right)x}^{\sqrt{4-x^{2}}}\!\!xdydx + \int_{1}^{2}\int_{-\sqrt{4-x^{2}}}^{\sqrt{4-x^{2}}}\!\!xdydx = \frac{4}{3}\sqrt{3}+\frac{5}{2}.
\end{align*}
}


Para considerar las coordenadas polares, se halla el ángulo de corte entre la recta inclinada y la circunferencia en el cuarto cuadrante; para ello, usaremos el punto
$(1,-\sqrt{3})$, entonces $\theta=\arctan \left(-\sqrt{3}\right)= -\frac{1}{3}\pi.$ La ecuación de la recta se reescribe como sigue:
$ r\sin \theta =1-\left( 1+\sqrt{3}\right)  r\cos \theta $,
de donde $r=\frac{1}{\sin(\theta)+\left(1+\sqrt{3}\right)\cos(\theta)}.$
Por lo tanto,
\[\ds \iint _R x dA=\int_{-\pi/3}^{\pi /2}\int_{\frac{1}{\sin(\theta)+\left( 1+\sqrt{3}\right) \cos (\theta) }}^{2}\left( r^{2}\cos \theta \right) drd\theta =\frac{5}{2}+\frac{4}{3}\sqrt{3}.\]
Las integrales anteriores se evalúan en \verb"Mathematica" mediante la ejecución de las siguientes instrucciones:




\begin{tcolorbox}[breakable,enhanced,title=\bf Instrucción en \verb"Mathematica", top=2pt,bottom=2pt,boxsep=0pt,before skip=2pt,after skip=0pt]
\begin{Verbatim}[formatcom=\letraverbatim]
Integrate[x,{y,-Sqrt[3],1},{x,(1-Sqrt[3])*(y-1)/2,
Sqrt[4-y^2]}]+Integrate[x,{y,1,2},{x,0,Sqrt[4-y^2]}]
Integrate[x,{x,0,1},{y,1-(1+Sqrt[3])*x,Sqrt[4-x^2]}]+
Integrate[x,{x,1,2},{y,-Sqrt[4-x^2],Sqrt[4-x^2]}]
Integrate[r^2*Cos[t],{t,-Pi/3,Pi/2},
{r,1/((1+Sqrt[3])*Cos[t]+Sin[t]),2}]
\end{Verbatim}
\end{tcolorbox}




%\vspace*{1.2em} 
\begin{enumerate}[start=9,leftmargin=0cm,itemindent=.75cm,labelwidth=\itemindent,labelsep=0cm,align=left,label={{\bf \arabic*.\,}}]
	\item 
Hallar el área de la región encerrada por la curva determinada por las ecuaciones paramétricas $x=2a\cos t-a\sin (2t),$ $y=b\sin t,$ $t\in \lbrack 0,2\pi ].$
\end{enumerate}


De acuerdo al teorema de Green, véase sección \secref{teoremadegreen}, el área $A$, de la región $R$ encerrada por la curva $C,$ puede hallarse por medio de la expresión
\[A=\iint_{R}dA=\int_{C}-\frac{1}{2}ydx+\frac{1}{2}xdy,\]
en este caso $x=2a\cos t-a\sin (2t),$ $y=b\sin (t).$ De estas ecuaciones se obtiene
$dx= -2a\sin t-2a\cos (2t),\,\,\, dy= b\cos t,$
entonces
\begin{multline*}
A=\frac{1}{2}\int_{0}^{2\pi}\bigg((2a\sin(t)+a\cos(2t))(b\sin(t))\\ +(2a\cos(t)-a\sin(2t))(b\cos t) \bigg) dt  = 2\pi ab.
\end{multline*}



\begin{figure}[H]
\centering
\caption{${\bf r}(t)=\langle 2a\cos t-a\sin (2t),b\sin t\rangle,$ $t\in [0,2\pi ]$}
\includegraphics[width=0.8\textwidth]{Fig/110.pdf}
\end{figure}

Otra posibilidad es la parametrización de la región interior a la curva.




\begin{tcolorbox}[breakable,enhanced,title=\bf Instrucción Mathematica, top=2pt,bottom=2pt,boxsep=0pt,
before skip=2pt,after skip=0pt]
Para $u\in [0,1]$, al multiplicar por $u$ las ecuaciones que determinan la curva, se define la siguiente función en $\mathbb{R}^3,$
${\bf r}(t,u)=\left\langle \left( 2a\cos t-\sin (2t) \right) u,ub\sin t,0\right\rangle.$
Con esta función se obtienen los siguientes vectores.
\begin{eqnarray*}
{\bf r}_{t} &=& \left\langle \left( -2a\sin (t)-2\cos(2t)
\right) u,ub\cos(t) ,0\right\rangle \\
{\bf r}_{u} &=& \left\langle 2a\cos (t) -\sin(2t) ,b\sin (t),0\right\rangle,
\end{eqnarray*}
luego, ${\bf r}_{t}\times {\bf r}_{u}= \left\langle 0,0,-ub\left( 2a+2\sin \left(t\right) \cos (2t) -\cos \left( t\right) \sin (2t)\right)  \right\rangle ,$
entonces
$\left\| {\bf r}_{t}\times {\bf r}_{u}\right\| =\left| bu\left( 2a+2\sin \left( t\right)
\cos (2t) -\cos \left( t\right) \sin (2t) \right)\right| ,$
cuya integral nos provee el área de la región en mención
\[\int_{0}^{2\pi }\int_{0}^{1}\left| bu\left( 2a+2\sin \left( t\right) \cos
(2t) -\cos \left( t\right) \sin (2t) \right) \right|
dudt= 2ab\pi.\]
\end{tcolorbox}



%
Para evaluar integrales dobles, es posible plantear un cambio de variable; la fundamentación teórica de este proceso se presentó en la sección \secref{cambioidobles}.
\section{
Cambio de variables en integrales dobles
}

\index{Cambio de variables}
Presentamos algunos ejemplos de aplicación del teorema de cambio de variables en integrales dobles.

\begin{enumerate}[leftmargin=0cm,itemindent=.75cm,labelwidth=\itemindent,labelsep=0cm,align=left,label={{\bf \arabic*.\,}}]
\item Evaluar $\ds \iint_{R}y^{2}dA,$ siendo $R$ la región del plano $xy$ acotada por las figuras de las ecuaciones.
$x=4-y^{2}$ y $x=4y^{2}-1.$


\begin{figure}[H]
	\centering	
	\caption{\small ${\bf T}(u,v)= (4u^2 - v^2, uv)$}
	\label{transformacion_uv}
	\begin{minipage}[b]{0.40\linewidth}
		\centering
		
		\includegraphics[width=\linewidth]{Fig/111.pdf}
	\end{minipage}
	\hspace{0.5cm} 
	\begin{minipage}[b]{0.35\linewidth}
		\centering
		
		\includegraphics[width=\linewidth]{Fig/112.pdf}
	\end{minipage}
	
	
	
	\fuentepropia
\end{figure}

Haremos la sustitución $x=4u^{2}-v^{2},\,\, y=uv.$
En la expresión $x=4-y^{2},$ se hace el cambio y se sigue que
$4u^{2}-v^{2}=4-\left( uv\right) ^{2}$ de donde $u^{2}\left( 4+v^{2}\right)
=4+v^{2},$ es decir,
$4+v^{2}=1$ (lo cual no es soluble en los números reales) o $u^{2}=1,$
entonces $u=\pm 1.$
Sustituyéndose en la ecuación $x=4y^{2}-1,$ se obtiene
$4u^{2}-v^{2}=4\left( uv\right) ^{2}-1,$ es decir, $4u^{2}\left(1-v^{2}\right) =v^{2}-1;$
por lo tanto $u=0$ o $v=\pm 1.$ De lo anterior se obtienen los límites para la región en el plano $uv.$

El jacobiano de la transformación es
\[\left| \begin{array}{cc}
\frac{\partial x}{\partial u} & \frac{\partial x}{\partial v} \\[0.2cm]
\frac{\partial y}{\partial u} & \frac{\partial y}{\partial u}
\end{array}
\right| =\left| \begin{array}{cc}8u & -2v \\[0.2cm]
v & u \end{array} \right| = 8u^{2}+2v^{2}.\]
Entonces,
\begin{eqnarray*}
\int_{-1}^{1}\int_{4y^{2}-1}^{4-y^{2}}y^{2}dxdy &=&
\int_{0}^{1}\int_{-1}^{1}\left( uv\right) ^{2}\left( 8u^{2}+2v^{2}\right) dudv \\
&=& \int_{0}^{1}\left( \frac{16}{5}v^{2}+\frac{4}{3}v^{4}\right), dv = \frac{4}{3}.
\end{eqnarray*}

\item Evaluar $\ds \iint_R xy dA,$ donde $R$ es la región del primer cuadrante acotada por las figuras de las expresiones
$y=1/x$, $y=3x$, $y=x$, $y=3/x$.




\begin{figure}[H]
\centering
\caption{${\bf T}(u,v)= (\sqrt{uv},\sqrt{\frac{u}{v}})$}
\includegraphics[width=0.8\textwidth]{Fig/113.pdf}
\fuentepropia
\end{figure}

Con las expresiones $y=1/x$ y $y=3/x,$ se puede hacer el cambio $xy=u,$ de esta manera $u\in \lbrack 1,3].$ Mientras que con $y=x$ y $y=3x,$ se puede hacer la sustitución $v=y/x,$ de modo que $v\in \lbrack 1,3].$
El jacobiano de la transformación es el recíproco del determinante de la matriz
\[\left| \begin{array}{cc}
\frac{\partial u}{\partial x} & \frac{\partial u}{\partial y} \\[0.1cm]
\frac{\partial v}{\partial x} & \frac{\partial v}{\partial y}
\end{array} \right| =\left| \begin{array}{cc}y & x \\ -y/x^{2} & 1/x\end{array}
\right| = \frac{2y}{x}=2v.\]
Entonces, el jacobiano de la transformación es $\frac{1}{2v};$ ahora bien,
{\footnotesize \[\iint_R xy dA=\int_{1}^{3}\int_{1}^{3}\frac{u}{2v}dudv= 2\ln 3.\]}
Esta integral también puede evaluarse en la forma
{\footnotesize
\begin{eqnarray*}
\iint_R xy dA &=& \int_{\frac{1}{\sqrt{3}}}^{1}\int_{\frac{1}{x}}^{3x}xydydx+\int_{1}^{\sqrt{3}}\int_{x}^{3/x}xydydx \\
&=& \int_{1}^{\sqrt{3}}\int_{\frac{y}{3}}^{y}xydxdy + \int_{\sqrt{3}}^3\int_{\frac{y}{3}}^{\frac{3}{y}}xydxdy =2\ln 3.
\end{eqnarray*}
}

Estas integrales se evalúan ejecutando las siguientes instrucciones:



\begin{tcolorbox}[breakable,enhanced,title=\bf Instrucción en \verb"Mathematica", top=2pt,bottom=2pt,boxsep=0pt,before skip=2pt,after skip=0pt]
\begin{Verbatim}[formatcom=\letraverbatim]
Integrate[x*y,{y,1,Sqrt[3]},{x,1/y,y}]+Integrate[x*y,
{y,Sqrt[3],3},{x,y/3,3/y}]
Integrate[x*y,{x,1/Sqrt[3],1},{y,1/x,3*x}]+
Integrate[x*y,{x,1,Sqrt[3]},{y,x,3/x}]
\end{Verbatim}
\end{tcolorbox}



\item Evaluar $\ds \iint_R xy dA,$ en donde $R$ es la región en el primer cuadrante acotada por las graficas de las expresiones
$x^{2}y=1,$ $x^{2}y=4,$ $y=x^{2}$ y $y=4x^{2}.$



\begin{figure}[H]
\centering
\caption{${\bf T}(u,v)= (\sqrt[4]{\frac{u}{v}},\sqrt{uv})$}
\includegraphics[width=0.8\textwidth]{Fig/114.pdf}
\end{figure}

Los puntos de intersección de estas curvas son $\left( 1,1\right) ,$ $\left( \sqrt{2},2\right) ,$
$\left( \frac{1}{\sqrt{2}},2\right) $ y $\left( 1,4\right)$; con esta información, se puede evaluar esta integral en coordenadas rectangulares, es decir

{\small
\begin{eqnarray*}
\ds \iint_R xy dA &=& \bi_{\frac{1}{\sqrt{2}}}^{1}\bi_{1/x^{2}}^{4x^{2}}xydydx+\bi_{1}^{\sqrt{2}}\bi_{x^{2}}^{4/x^{2}}xydydx \\
&=& \bi_{1}^{2}\bi_{\sqrt{\frac{1}{y}}}^{\sqrt{y}}xydxdy+\bi_{2}^{4}\bi_{\sqrt{\frac{y}{4}}}^{\sqrt{\frac{4}{y}}}xydxdy =\frac{7}{3}
\end{eqnarray*}
}
Estas integrales se evalúan con las siguientes instrucciones:



\begin{tcolorbox}[breakable,enhanced,title=\bf Instrucción en \verb"Mathematica", top=2pt,bottom=2pt,boxsep=0pt,before skip=2pt,after skip=0pt]
\begin{Verbatim}[formatcom=\letraverbatim]
Integrate[x*y,{y,2,4},{x,Sqrt[y/4],Sqrt[4/y]}]+
Integrate[x*y,{y,1,2},{x,Sqrt[1/y],Sqrt[y]}]
		
Integrate[x*y,{x,1/Sqrt[2],1},{y,1/x^2,4*x^2}]+
Integrate[x*y,{x,1,Sqrt[2]},{y,x^2,4/x^2}]
\end{Verbatim}
\end{tcolorbox}




También se puede utilizar el teorema de cambio de variable, definiendo $u=x^{2}y$ y $v=y/x^{2},$ de donde $x=\sqrt[4]{\frac{u}{v}},\,\, y=\sqrt{u}\sqrt{v}$,
por lo tanto, el jacobiano de la transformación es
{\small \[\left| \begin{array}{cc} \frac{\partial }{\partial u}\left( \sqrt[4]{\frac{u}{v}}\right) & \frac{\partial }{\partial v}\left( \sqrt[4]{\frac{u}{v}}\right) \\[0.1cm]
\frac{\partial }{\partial u}\left( \sqrt{uv}\right) & \frac{\partial }{\partial v}\left( \sqrt{uv}\right) \end{array} \right| = \frac{1}{4\sqrt[4]{uv^{3}}}.\]}
A partir de la definición, es posible establecer que $1\leq u \leq 4$ y $1\leq v \leq 4$, respectivamente.
{\small
\begin{multline*}
\iint_R xy dA = \int_{1}^{4}\int_{1}^{4}\left( \sqrt[4]{\frac{u}{v}}\right) \left( \sqrt{u}\sqrt{v}\right) \frac{1}{4\sqrt[4]{uv^{3}}}\, dudv\\
=\int_{1}^{4}\int_{1}^{4}\frac{1}{4}\frac{\sqrt{u}}{\sqrt{v}}\, dudv= \frac{7}{3}.
\end{multline*}
}

\item
Evaluar la integral {\footnotesize $\ds\int_{0}^{1}\int_{2y}^{\sqrt{1+3y^2}}xy^{2}dxdy+\int_{-1}^{0}\int_{-2y}^{\sqrt{1+3y^2}}xy^{2}dxdy.$} La región de integración está limitada por las figuras de las ecuaciones $x= \pm 2y\,$ y
$\, x^2-3y^2=1,$ con $x\geq 0$.



\begin{figure}[H]
\centering
\caption{${\bf T}(u,v)= (u\cosh v, \frac{u}{\sqrt{3}}\sinh v)$}
\includegraphics[width=0.4\textwidth]{Fig/115.pdf} \hspace{5mm}
\includegraphics[width=0.4\textwidth]{Fig/116.pdf}
\end{figure}

Puede probarse que
\[ \int_{0}^{1}\int_{2y}^{\sqrt{1+3y^{2}}}xy^{2}dxdy+\int_{-1}^{0}\int_{-2y}^{\sqrt{1+3y^{2}}}xy^{2}dxdy= \frac{2}{15}.\]
El siguiente proceso permite evaluar las anteriores integrales de una manera distinta.
 
Al hacer $x=u\cosh v,$ $y=\frac{u}{\sqrt{3}}\sinh v,$ se sigue que la expresión $x^{2}-3y^{2}=1$ se transforma en
$\left( u\cosh v\right) ^{2}-3(\frac{u}{\sqrt{3}}\sinh v)^2=1,$
que se reduce a $u^{2}=1,$ es decir, $u=\pm 1.$
Para la recta $y=x/2,$ se sigue que
$2\left( \frac{u}{\sqrt{3}}\right) \sinh v=u\cosh v$, es decir $v=\ln \left(
2+\sqrt{3}\right) ,$ mientras que para la recta $y=-x/2,$ se concluye que
$v=-\ln \left(2+\sqrt{3}\right). $
  
Con el cambio propuesto, el jacobiano de la transformación es
%
\begin{eqnarray*}
\left| \begin{array}{ll}
\frac{\partial }{\partial u}\left( u\cosh v\right) & \frac{\partial }{\partial v}\left( u\cosh v\right)  \\[0.2cm] \frac{\partial }{\partial u}\left( \frac{u}{\sqrt{3}}\sinh v\right)  & \frac{
\partial }{\partial v}\left( \frac{u}{\sqrt{3}}\sinh v\right)
\end{array}\right| &=& \frac{\sqrt{3}}{3}u\cosh ^{2}\left( v\right) -\frac{\sqrt{3}}{3}u\sinh ^{2}\left( v\right) \\&=&  \frac{1}{3}u\sqrt{3}.
\end{eqnarray*}
Entonces, la integral en mención se transforma en
\[
\int_{0}^{1} \int_{-\ln \left( 2+\sqrt{3} \right)}^{\ln \left( 2+\sqrt{3} \right)}
\left( u \cosh v \right)
\left( \frac{u}{\sqrt{3}} \sinh v \right)^2
\left( \frac{\sqrt{3}}{3} u \right)
\, dv \, du = \frac{2}{15}.
\]

%

\begin{tcolorbox}[breakable,enhanced,title=\bf Instrucción en \verb"Mathematica", top=2pt,bottom=2pt,boxsep=0pt,before skip=2pt,after skip=0pt]
\begin{Verbatim}[formatcom=\letraverbatim]
Integrate[x*y^2,{y,0,1},{x,2*y,Sqrt[1+3*y^2]}]+
Integrate[x*y^2,{y,-1,0},{x,-2*y,Sqrt[1+3*y^2]}]

Integrate[u*Cosh[v]*(u/Sqrt[3]/3*u),
{u,0,1},{v,-Log[2+Sqrt[3]],Log[2+Sqrt[3]]}]
\end{Verbatim}
\end{tcolorbox}


 

\item Sea $R$ el trapecio con vértices $\left( 0,2\right),$ $\left( 1,1\right),\left( 2,2\right),$ $\left( 1,3\right).$
En coordenadas rectangulares, la integral $\ds \iint _{R}\frac{x-y}{x+y}dA$ se puede evaluar como se muestra a continuación:
\[I=\int_{0}^{1}\int_{2-x}^{2+x}\frac{x-y}{x+y}dydx+\int_{1}^2\int_{x}^{4-x}\frac{x-y}{x+y}dydx=-\ln \left( 2\right).\]
Al hacer el cambio $u=x+y,$ $v=y-x,$ se sigue que $y=\frac{u+v}{2},$ $x=\frac{u-v}{2}$. 



\begin{figure}[H]
\centering
\caption{\small ${\bf T}(u,v)= \left(\frac{u - v}{2},\, \frac{u + v}{2}\right)$}
\includegraphics[width=0.4\textwidth]{Fig/117.pdf} \hspace{5mm}
\includegraphics[width=0.4\textwidth]{Fig/118.pdf}
\end{figure}

La figura muestra la región transformada; el jacobiano de la transformación es
$J=\det \left( \begin{array}{cc}
\frac{1}{2} & -\frac{1}{2} \\[0.2cm]
\frac{1}{2} & \frac{1}{2} \end{array} \right) = \frac{1}{2}$; luego
$I=\ds \int_{2}^{4}\ds\int_{0}^2\frac{-v}{u}\left( \frac{1}{2}\right) dvdu= -\ln \left( 2\right) .$

\item Evaluar $\ds \iint_R \sqrt{xy}\ln \left( \frac{y}{x}\right) dA,$ en donde $R$ es la región en el plano $xy$ determinada por el cambio de variables
$u^{2}=xy,$ $v=\ln \left( y/x\right) ,$ con $u\in \lbrack 1,2],$ $v\in \lbrack 0,2].$ \\[0.2cm]

Con el cambio de variable se sigue que $e^{v}=y/x,$: es decir, $y=xe^{v},$
que, al reemplazar en la primera ecuación, se obtiene que
$u^{2}=x^{2}e^{v},$ entonces $x=ue^{-v/2},$ y $y=ue^{v/2}.$
Para determinar las ecuaciones de las curvas que acotan la región, lo haremos por casos.
\begin{itemize}[label=$\bullet$\,\,]
\item Si $u=1$ y $0\leq v\leq 2,$
entonces, $x=e^{-v/2}$ y $y=e^{v/2},$ por tanto $xy=1.$
\item Si $u=2$ y $0\leq v\leq 2,$ entonces \
$x=2e^{-v/2}$ y $y=2e^{v/2}$ es decir $xy=4.$
\item Si $v=0$ y $1\leq u\leq 2,$
se sigue que $x=u$ y $y=u$ entonces $x=y.$
\item Si $v=2$ y $1\leq u\leq 2$
Entonces, $x=ue^{-1}$ y $y=ue$; por lo tanto, $y=e^{2}x.$
\end{itemize}

El siguiente gráfico muestra la región de integración. Los puntos de corte de las expresiones que determinan dicho conjunto en el plano $xy$
son $(e^{-1},e),$ $(1,1)$, $(2e^{-1},2e)$ y $(2,2).$ También se muestra la región en el plano $uv,$
que se obtiene al aplicar la transformación.

\begin{figure}[H]
	\centering
	\caption{${\bf T}(u,v)= (ue^{-v/2},ue^{v/2})$}
	\begin{minipage}{0.37\linewidth}
		\centering
		\includegraphics[width=\linewidth]{Fig/119.pdf}
	\end{minipage}
	\hspace{1em}
	\begin{minipage}{0.36\linewidth}
		\centering
		\includegraphics[width=\linewidth]{Fig/120.pdf}
	\end{minipage}
	
	
	\fuentepropia	
\end{figure}

La integral en mención puede evaluarse usando coordenadas rectangulares como sigue
\begin{multline*}
\int_{e^{-1}}^{1}\int_{1/x}^{e^{2}x}\sqrt{xy}\ln \left( \frac{y}{x}\right)
dydx+\int_{1}^{2e^{-1}}\int_{x}^{e^{2}x}\sqrt{xy}\ln \left( \frac{y}{x}\right) dydx\\
+\int_{2e^{-1}}^{2}\int_{x}^{4/x}\sqrt{xy}\ln \left( \frac{y}{x}\right) dydx= \frac{14}{3}
\end{multline*}
Esta integral se evalúa utilizando las siguientes instrucciones.



\begin{tcolorbox}[breakable,enhanced,title=\bf Instrucción en \verb"Mathematica", top=2pt,bottom=2pt,boxsep=0pt,before skip=2pt,after skip=0pt]
\begin{Verbatim}[formatcom=\letraverbatim]
Integrate[Sqrt[x*y]*Log[y/x],{x,Exp[-1],1},
{y,1/x,Exp[2]*x}]+
Integrate[Sqrt[x*y]*Log[y/x],{x,1,2*Exp[-1]},
{y,x,Exp[2]*x}]+
Integrate[Sqrt[x*y]*Log[y/x],{x,2*Exp[-1],2},
{y,x,4/x}]
\end{Verbatim}
\end{tcolorbox}



Utilizando el cambio de variables, el jacobiano de la transformación es
\[
\det \left(
\begin{array}{cc}
\frac{\partial }{\partial u}\left( ue^{-v/2}\right)  & \frac{\partial }{\partial v}\left( ue^{-v/2}\right)  \\[0.2cm]
\frac{\partial }{\partial u}\left( ue^{v/2}\right)  & \frac{\partial }{\partial v}\left( ue^{v/2}\right)
\end{array}
\right)
=
\det \left(
\begin{array}{cc}
e^{-\frac{1}{2}v} & -\frac{1}{2}ue^{-\frac{1}{2}v} \\
e^{\frac{1}{2}v} & \frac{1}{2}ue^{\frac{1}{2}v}
\end{array}
\right)
= u
\]

Con lo cual, la integral se reduce a
\[\iint_R \sqrt{xy}\ln \left( \frac{y}{x}\right) dA= \int_{1}^{2}\int_{0}^{2}u^{2}vdvdu= \frac{14}{3}\]

\item Aplicar el cambio de variables $u=\sqrt{xy},\,\, v=\sqrt{\frac{y}{x}},$ en la integral
\[\int_{1}^{4} \int_{1/y}^{y}\sqrt{\frac{y}{x}}e^{\sqrt{xy}}dxdy.\]
Del cambio se sigue que
\[uv=\sqrt{xy}\sqrt{\frac{y}{x}}= y;\] además, \[x=\frac{u^{2}}{uv}= \frac{u}{v}.\]
El jacobiano en este caso es
\[\left|  \begin{array}{cc}
\frac{1}{v} & v \\[0.1cm]
\frac{-u}{v^{2}} & u
\end{array} \right| = \frac{2u}{v}.\]
Las respectivas regiones en los planos $xy$ y $uv$
se muestran enseguida.

\begin{figure}[h!]
	\centering
	\caption{${\bf T}(u,v)= \left(\frac{u}{v},\, uv\right)$}
	\begin{minipage}{0.35\linewidth}
		\centering
		\includegraphics[width=\linewidth]{Fig/121.pdf}
	\end{minipage}
	\hspace{1em}
	\begin{minipage}{0.40\linewidth}
		\centering
		\includegraphics[width=\linewidth]{Fig/122.pdf}
	\end{minipage}
	
	
	\fuentepropia
\end{figure}

Por ello, la integral se transforma en
\[\int_{1}^{4}\int_{1/y}^{y}\sqrt{\frac{y}{x}}e^{\sqrt{xy}}dxdy
=\int_{1}^{4}\int_{1}^{4/u}2ue^{u}dvdu = 2e^4-8e\]
%\begin{eqnarray*} \int_{1}^{4}\int_{1/y}^{y}\sqrt{\frac{y}{x}}e^{\sqrt{xy}}dxdy &=&
%\int_{1}^{4}\int_{1}^{4/u}ve^{u}\left( \frac{2u}{v}\right)dvdu \\
%&=&\int_{1}^{4}\int_{1}^{4%/u}2ue^{u}dvdu = 2e^4-8e %\end{eqnarray*}

Estas integrales se pueden evaluar con ayuda de las siguientes sentencias.

\vspace{3pt}

\begin{tcolorbox}[breakable,enhanced,title=\bf Instrucción en \verb"Mathematica", top=2pt,bottom=2pt,boxsep=0pt,before skip=2pt,after skip=0pt]
\begin{Verbatim}[formatcom=\letraverbatim]
Integrate[Sqrt[y/x]*Exp[Sqrt[x*y]],{y,1,4},{x,1/y,y}]
Integrate[2*u*Exp[u],{u,1,4},{v,1,4/u}]
\end{Verbatim}
\end{tcolorbox}

\end{enumerate}

\begin{ejer}
Use el teorema de cambio de variable para resolver los siguientes problemas

\begin{enumerate}
\item   Evaluar
$\displaystyle \iint_R (x^2+y^2)dA, $ donde $R$ es la región en el primer cuadrante acotada por las figuras de las expresiones $x^{2}-y^{2}=3,$ $y^{2}-x^{2}=3,$ $xy=2,$ $xy=3\sqrt{2}.$
Muestre también que el área de la región mencionada está dada por $ -6+9\sqrt{2}.$



\begin{figure}[H]
\centering
\caption{${\bf T}(u,v)= (x,y)$}
\includegraphics[width=0.8\textwidth]{Fig/123.pdf}
\end{figure}

\begin{comment}
$ \int_{1}^{\sqrt{3}}\int_{2/x}^{\sqrt{3+x^{2}}}(x^2+y^2)dydx+\int_{\sqrt{3}}^{2}\int_{2/x}^{3\sqrt{2}/x}(x^2+y^2)dydx+\int_{2}^{\sqrt{6}}\int_{\sqrt{x^{2}-3}}^{3\sqrt{2}/x}(x^2+y^2)dydx=-6+9\sqrt{2}$

Es claro que $J= \frac{1}{2(x^2+y^2)}.$
Con lo cual $\int_{-3}^{3}\int_{2}^{3\sqrt{2}} \frac{1}{2} dudv= -6+9\sqrt{2} $
\end{comment}

\item Determine el área de la región $R$ acotada en el primer cuadrante para las curvas
$x^2+y^2=2$, $x^2+y^2=4,$ $x^2-y^2=2$, $x^2+y^2=4$.



\begin{figure}[H]
\centering
\caption{${\bf T}(u,v)= (x,y)$}
\includegraphics[width=0.8\textwidth]{Fig/124.pdf}
\end{figure}

\end{enumerate}
\end{ejer}

\section{Áreas de superficies parametrizadas} \index{Superficies parametrizadas!áreas}
Además de parametrizar algunas superficies, hallaremos su área. Esta tarea se hará básicamente utilizando la expresión
\ref{areasup}.
\subsection*{Un cono truncado oblicuo}
Las figuras de las funciones vectoriales $\overrightarrow{\alpha}(t)=\langle 1+\cos(t) ,\sin(t),2\rangle $ y
$\overrightarrow{\beta}(t)=\langle 3\cos (t),3\sin(t),0\rangle ,$ para $t\in [0,2\pi],$ son dos circunferencias.
Formando una combinación convexa con un punto $P$ sobre la curva $\overrightarrow{\alpha}$ y un punto $Q$ sobre la curva $\overrightarrow{\beta},$ se genera un cono truncado oblicuo.


\begin{figure}[H]
\centering
\caption{Cono truncado oblicuo determinado por ${\bf r}(u,t)$}
\includegraphics[width=0.4\textwidth]{Fig/125.pdf} \hspace{5mm}
\includegraphics[width=0.4\textwidth]{Fig/126.pdf}
\fuentepropia
\end{figure}

Una parametrización para esta región es la siguiente.
\begin{equation*}
\resizebox{0.97\textwidth}{!}{$
\begin{aligned}
{\bf r}(u,t) &= (1-u) \overrightarrow{\alpha}(t) + u\overrightarrow{\beta}(t) \\
&= \langle 1 - u + (1 + 2u)\cos t,\ (1 + 2u)\sin t,\ 2 - 2u \rangle,\!\!\! u \in [0,1],\ t \in [0,2\pi].
\end{aligned}
$}
\end{equation*}
%$u\in[0,1],$ %$t\in[0,2\pi].$
Para calcular el área superficial, se requieren los siguientes vectores
{\small
\begin{eqnarray*}
{\bf r}_{t} &=& \langle -\sin t( 1+2u) ,( 1+2u) \cos t,0\rangle \\
{\bf r}_{u} &=& \langle -1+2\cos t,2\sin t,-2\rangle. \end{eqnarray*}
}
además,
$ {\bf r}_{t}\times {\bf r}_{u}= \langle -2( 1+2u) \cos t,-2( \sin t) (1+2u) ,( \cos t-2) ( 1+2u) \rangle .$
El elemento diferencial de área es
$\| {\bf r}_{t}\times {\bf r}_{u}\| = (1+2u) \sqrt{\cos ^{2}t-4\cos t+8}.$
Por último, el área superficial está dada por
{\small
\[\int_{0}^{1}\int_{0}^{2\pi }( 1+2u) \sqrt{\cos ^{2}t-4\cos t+8}\, dtdu\approx 36.117.\]}
Este valor se consiguió utilizando las siguientes instrucciones.



\begin{tcolorbox}[breakable,enhanced,title=\bf Instrucción en \verb"Mathematica", top=2pt,bottom=2pt,boxsep=0pt,before skip=2pt,after skip=0pt]
\begin{Verbatim}[formatcom=\letraverbatim]
r=(1-u)*{1+Cos[t],Sin[t],2}+u*{3*Cos[t],3*Sin[t],0};
A=Assuming[t>=0&&u>=0,Simplify[Norm[Cross[D[r,u],
D[r,t]]]]]
NIntegrate[A,{t,0,2*Pi},{u,0,1}]
\end{Verbatim}
\end{tcolorbox}



\index{Cilindro cardiodico}
\subsection*{Una esfera y dos cilindros perpendiculares}
La esfera centrada en el origen de radio $2$ y los cilindros con ecuaciones son $x^{2}+y^{2}=2x$ y $x^{2}+z^{2}=2x$
se intersecan. Se determinará la superficie sobre la esfera acotada por los cilindros, y de esta nos interesará su área. % Los siguientes gráficos ilustran la situación mencionada.


\begin{figure}[H]
\centering
\caption{Esfera y dos cilindros}
\includegraphics[width=0.4\textwidth]{Img/6g.pdf} \hspace{5mm}
\includegraphics[width=0.4\textwidth]{Img/6a.pdf}

\includegraphics[width=0.4\textwidth]{Img/6b.pdf}
\fuentepropia
\end{figure}

La porción sobre la esfera interior de los cilindros se parametriza en dos casos:

\newlength{\bulletindent} 
\settowidth{\bulletindent}{\textcolor{black}{\scriptsize\ding{110}}\quad}
\begin{minipage}[t]{\textwidth}
 
\begin{itemize}[
leftmargin=\dimexpr 1.5cm - \bulletindent\relax,itemsep=0pt,parsep=0pt,labelsep=0.4em,label=\textcolor{black}{\scriptsize\ding{110}}]
\item El cilindro $x^2+y^2=2r$ se reescribe como $(x-1)^2+y^2=1$. Para la superficie sobre la esfera se puede hacer $x=1+u\cos(t)$, $y=u\sin(t)$, donde $t\in[0,2\pi]$, lo cual parametriza la región sobre la esfera interior al cilindro proyectada el plano $xy$. Reemplazando en la ecuación de la esfera se obtiene {\small $-3+2u\cos(t)+u^2+z^2=0$}, cuya solución es {\small $z=\pm\sqrt{3-2u\cos(t)-u^2}$}. Por lo tanto la función
\end{itemize}
\end{minipage}

{\footnotesize
\[{\bf r}( u,t) =\left\langle 1+u\cos (t) ,u\sin \left(
t\right) ,\sqrt{3-2u\cos (t) -u^{2}}\right\rangle,\]}
con $t\in \lbrack 0,2\pi ],\,\, u\in \lbrack 0,1],$
representa la parte de la esfera que se halla acotada por el cilindro, la otra parte se determina con la raíz negativa en la tercera componente.

%\clearpage


\begin{figure}[H]
\centering
\caption{Curva de intersección de la esfera y los cilindros}
\includegraphics[width=0.4\textwidth]{Img/6f.pdf} \hspace{5mm}
\includegraphics[width=0.4\textwidth]{Img/6c.pdf}
\fuentepropia
\end{figure}

\settowidth{\bulletindent}{\textcolor{black}{\scriptsize\ding{110}}\quad}
\begin{minipage}[t]{\textwidth}


\begin{itemize}[
leftmargin=\dimexpr 1.5cm - \bulletindent\relax,itemsep=0pt,parsep=0pt,labelsep=0.4em,before=\vspace{-0.4cm},label=\textcolor{black}{\scriptsize\ding{110}}]
\item Para la expresión $x^{2}+z^{2}=2x$ (que representa el otro cilindro), el tratamiento es similar al caso anterior, obteniendo la función

{\footnotesize
\[
{\\bf r}_1(u,t) = \left\langle 1 + u\cos(t),\, \pm \sqrt{3 - 2u\cos(t) - u^{2}},\, u\sin(t) \right\rangle
\]
}

con $t \in [0, 2\pi]$, $u \in [0, 1]$, la cual parametriza la porción de la esfera dentro del cilindro con ecuación $x^{2}+z^{2}=2x$.
\end{itemize}
\end{minipage}

%


\begin{figure}[H]
\centering
\caption{Curva de intersección de la esfera y los cilindros}
\includegraphics[width=0.4\textwidth]{Img/6d.pdf} \hspace{5mm}
\includegraphics[width=0.4\textwidth]{Img/6e.pdf}

\includegraphics[width=0.4\textwidth]{Fig/127.pdf}
\fuentepropia
\end{figure}

Los anteriores gráficos se obtienen ejecutando las siguientes instrucciones:


\begin{enumerate}[leftmargin=0cm,itemindent=.75cm,labelwidth=\itemindent,labelsep=0cm,align=left,label={{\bf \arabic*.\,}}]
\item La esfera y los cilindros

Región sobre la esfera acotada por los cilindros, se definen tres objetos nombrados \verb"a,b,c", luego se muestran en un solo gráfico.



\begin{tcolorbox}[breakable,enhanced,title=\bf Instrucción en \verb"Mathematica", top=2pt,bottom=2pt,boxsep=0pt,before skip=2pt,after skip=0pt]
\begin{Verbatim}[formatcom=\letraverbatim]
ContourPlot3D[{x^2+y^2+z^2==4,x^2+y^2==2*x,x^2+
z^2==2*x},{x,-2,2},{y,-2,2},{z,-2,2},PlotPoints->80,
Mesh->False]
ContourStyle->{Orange,Yellow,LightGreen}
\end{Verbatim}
\end{tcolorbox}




\item Región sobre la esfera acotada por los cilindros, se definen tres objetos nombrados \verb"a,b,c", luego se muestran en una sola figura usando la función \textit{show}.



\begin{tcolorbox}[breakable,enhanced,title=\bf Instrucción en \verb"Mathematica", top=2pt,bottom=2pt,boxsep=0pt,before skip=2pt,after skip=0pt]
\begin{Verbatim}[formatcom=\letraverbatim]
Clear[x,y,z]
a=ContourPlot3D[x^2+y^2+z^2==3.95,{x,-2,2},{y,-2,2},
{z,-2,2},ContourStyle->{Cyan,Opacity[0.4]},Mesh->Fals
PlotPoints->50];
r1={x,y,z}={1+Cos[t],Sqrt[2-2*Cos[t]],Sin[t]};
r5={x1,y1,z1}={1+u*Cos[t],Sqrt[3-2*u*Cos[t]-u^2],
u*Sin[t]};
b=ParametricPlot3D[{r1,{x,-y,z},{x,z,y},{x,z,-y}},
{t,0,2*Pi},PlotStyle->Black,PlotPoints->80];
c=ParametricPlot3D[{r5,{x1,-y1,z1},{x1,z1,y1},
{x1,z1,-y1}},{t,0,2*Pi},{u,0,1},
PlotStyle->Blend[{Yellow,Pink,Green}]];Mesh->False,
Show[a,b,c]
\end{Verbatim}
\end{tcolorbox}



Por la simetría de la región en mención, el área de la porción sobre la esfera interior a los cilindros corresponde a cuatro veces la integral del término
$\left\Vert {\bf r}_{t}\times {\bf r}_{u}\right\Vert,$ que es el elemento diferencial del área.
Para tal fin, se hallan las siguientes funciones vectoriales:

\begin{eqnarray*}
{\bf r}_{t}(t,u) &=& \la -u\sin t,u\cos t,\frac{u\sin t}{\sqrt{3-u^2 -2u\cos t}}\ra \\
{\bf r}_{u}(t,u)&=& \la\cos t,\sin t,-\frac{u+\cos t}{\sqrt{3-u^2-2u\cos t}}\ra \\
{\bf r}_{t}\times{\bf r}_{u}&=&\la-\frac{u^2 \cos t+u}{\sqrt{3-u^2 -2u\cos t}},\frac{-u^2 \sin t}{\sqrt{3-u^2 -2u\cos t}},-u,\ra,
\end{eqnarray*}
\end{enumerate}
entonces
$\left\Vert {\bf r}_{t}\times{\bf r}_{u} \right\Vert =\frac{2u}{\sqrt{3-u^{2}-2u\cos t}}.$
Por lo anterior, el área $A$ de la superficie es cuatro veces la integral del término anterior; es decir,

\[4\int_{0}^{2\pi}\int_{0}^{1}\frac{2u}{\sqrt{3-u^{2}-2u\cos t}}dudt \approx 18. 265.\]

Este cálculo se hace con \verb"Mathematica" ejecutando las siguientes instrucciones:



\begin{tcolorbox}[breakable,enhanced,title=\bf Instrucción en \verb"Mathematica", top=2pt,bottom=2pt,boxsep=0pt,before skip=2pt,after skip=0pt]
\begin{Verbatim}[formatcom=\letraverbatim]
{x,y,z}={1+u*Cos[t],Sqrt}[3-u^2-2*u*Cos[t]]};
r={x,y,z};
A=Norm[Cross[D[r,t],D[r,u]]];
NIntegrate[4*A,{u,0,1},{t,0,2Pi}]
\end{Verbatim}
\end{tcolorbox}



\subsection*{Un cilindro cardiogénico bajo un paraboloide}

Las ecuaciones $3z=9-(x+2)^2-(y+1)^2$ y $r=1-\cos(t)$ representan un paraboloide y un cilindro cardiodónico, respectivamente.
Estas superficies, junto con el plano $z=0$, acotan un sólido.
En el plano $xy$ el cilindro se proyecta como un cardioide que puede parametrizarse como
\[\overrightarrow{s}(t)=\la(1-\cos(t))\cos(t),(1-\cos(t))\sin(t),0\ra, \,\, t \in [0,2\pi].\]


Sobre el paraboloide, la coordenada $z$ tiene la forma

{\small
\begin{align*}
z &= \frac{1}{3}\left(9 - \left((1 - \cos(t))\cos(t) + 2\right)^2 - \left((1 - \cos(t))\sin(t) + 1\right)^2\right) \\
&= \frac{1}{6}(13 - 4\sin(t) + 2\sin(2t) + \cos(2t))
\end{align*}
}




\begin{figure}[H]
\centering
\caption{\small Cilindro cardiogénico y paraboloide}
\includegraphics[width=0.4\textwidth]{Img/40a.pdf} \hspace{5mm}
\includegraphics[width=0.4\textwidth]{Img/40b.pdf}
\fuentepropia
\end{figure}

Por tanto, la curva de intersección se parametriza por medio de la función

{\footnotesize
\begin{multline*} \overrightarrow{r}(t)=\big \langle(1-\cos(t))\cos(t),(1-\cos(t))\sin(t), \\ \frac{1}{3}\left(13-4\sin(t)+2\sin(2t)+\cos(2t)\right)\bigg \rangle,\,\, t \in [0,2\pi].
\end{multline*}
}



\begin{figure}[H]
\centering
\caption{Superficie sobre el cilindro cardiogénico acotada por el paraboloide}
\includegraphics[width=0.4\textwidth]{Fig/128.pdf} \hspace{5mm}
\includegraphics[width=0.4\textwidth]{Img/39c.pdf}
\fuentepropia
\end{figure}

Una combinación convexa de las curvas $\overrightarrow{s}$ y $\overrightarrow{r}$ de la forma $(1-u)\overrightarrow{s}(t)+u\overrightarrow{r}(t),$
parametriza la porción del cilindro que se hallaba bajo el paraboloide y sobre el plano $xy.$
{\small
\begin{multline*} \overrightarrow{p}(t,u)=\big\langle (1-\cos(t))\cos(t),(1-\cos(t))\sin(t), \\ \frac{1}{6}\left(u(13-4\sin(t)+2\sin(2t)+\cos(2t)) \right)\big\rangle,
\end{multline*}
}

Calcularemos el área de esta región. Con esta función es sencillo ver que
\begin{eqnarray*}
\overrightarrow{p}_{t}&=& \bigg \langle-\sin(t)+2\sin(t)\cos(t),\cos(t)-2\cos^{2}(t)+1,\\
&& \hspace{3.8cm} \frac{1}{6} u (-2 \sin (2 t)-4 \cos (t)+4\cos(2t)) \bigg \rangle \\
\overrightarrow{p}_{u}&=& \la 0,0,\frac{1}{6} (-4 \sin (t)+2 \sin (2 t)+\cos (2 t)+13) \ra
\end{eqnarray*}
Con estos vectores obtenemos el siguiente término:
\begin{multline*}
\|\overrightarrow{p}_{t}\times \overrightarrow{p}_{u}\|  =
\frac{1}{6}\bigg |4\cos\left(\frac{t}{2}\right)-6\cos\left(\frac{3t}{2}\right)+2\cos\left(\frac{5t}{2}\right)-22\sin\left(\frac{t}{2}\right)\\
+7\sin\left(\frac{3t}{2}\right)-3\sin\left(\frac{5t}{2}\right)\bigg |
\end{multline*}
Al integrarlo con los límites descritos antes, proporciona el valor del área del cilindro cardiogénico que se encuentra entre el plano $xy$ y el paraboloide, esto es
\[\int_{0}^{2\pi}\int_{0}^{1}\|\overrightarrow{p}_{t}(t,u)\times \overrightarrow{p}_{u}(t,u)\|dudt=\frac{608}{45}.\]
Los cálculos se hacen en \verb"Mathematica", para este fin se define $x$ e $y$, que corresponden a la parametrización de la cardioide en el plano. Sustituyendo esto en la ecuación del paraboloide, se define
$z$, que al unir los puntos permite parametrizar el sólido mediante la expresión $r$.
Luego se obtiene el gráfico de $r$ y se calcula el producto triple que se denomina $A$, este se integra con los límites descritos anteriormente.




\begin{tcolorbox}[breakable,enhanced,title=\bf Instrucción en \verb"Mathematica", top=2pt,bottom=2pt,boxsep=0pt,before skip=2pt,after skip=0pt]
\begin{Verbatim}[formatcom=\letraverbatim]
{x,y}={(1-Cos[t])*Cos[t],(1-Cos[t])*Sin[t]};
z=(9-(x+2)^2-(y+1)^2)/3;
r=Simplify[(1-u)*{x,y,z}+u*{x,y,0}];
ParametricPlot3D[r,{t,0,2*Pi},{u,0,1}];
A=Assuming[t>=0,Simplify[Norm[Cross[D[r,t],D[r,u]]]]]
Integrate[Abs[A],{t,0,2*Pi},{u,0,1}]
\end{Verbatim}
\end{tcolorbox}




\subsection*{Un paraboloide fuera de un cilindro cardiogénico}
Consideremos el paraboloide con ecuación\ $3z=9-x^{2}-y^{2}$ y el cilindro cardiodónico $r=1-\cos(t)$.
La proyección del paraboloide en el plano $xy$ es la circunferencia centrada en el origen del radio $3$,
que puede parametrizarse como $\la 3\cos(t),3\sin(t)\ra $.
Para $t\in[0,2\pi],\, u\in[0,1]$, la región que se encuentra entre la cardioide y la circunferencia se parametriza por:

{\small
\begin{multline*}
(1-u)\la(1-\cos(t))\cos(t),(1-\cos(t))\sin(t)\ra+u\la 3\cos(t),3\sin(t)\ra \\
= \la(1+2u+(u-1)\cos(t))\cos(t),(1+2u+(u-1)\cos(t))\sin(t)\ra.
\end{multline*}
}



\begin{figure}[H]
\centering
\caption{Superficie sobre el paraboloide $3z=9-x^{2}-y^{2}$}
\includegraphics[width=0.4\textwidth]{Fig/129.pdf} \hspace{5mm}
\includegraphics[width=0.4\textwidth]{Fig/130.pdf}
\fuentepropia
\end{figure}

\par 

La coordenada $z$ se obtiene reemplazando $x$ y $z$ en la ecuación del paraboloide, esto es

{\small
\begin{align*}
	z &= \frac{1}{3}(9-x^2-y^2) \\
	&= \frac{1}{3}\left(9 - \left((1+2u + (u-1)\cos(t))\cos(t)\right)^2 \right. \\
	&\quad \left. - \left((1+2u + (u-1)\cos(t))\sin(t)\right)^2\right) \\
	&= \frac{1}{3}\left( (u-1)\cos(t)+2 \right) \left( (1-u)\cos(t) - 2u - 4 \right).
\end{align*}
}
Entonces la superficie sobre el paraboloide que se encuentra fuera del cilindro cardiogénico, y sobre el plano $xy$ se parametriza como:
{\small
\begin{multline*}
{\bf r}(t,u)=\bigg \langle(1+2u+(u-1)\cos(t))\cos(t),(1+2u+(u-1)\cos(t))\sin(t) \\
	\frac{1}{3}(u-1)(\cos(t)+2)((1-u)\cos(t)-2u-4) \bigg \rangle, u\in\lbrack0,1],t\in [0,2\pi].
\end{multline*}
}
%$((\cos(t))(1+2u+(u-1)\cos(t)),(1+2u+(u-1)\cos(t))\sin(t),\frac{1}{3}(u-1)(\cos(t)+2)((1-u)\cos(t)-2u-4))$

\begin{figure}[H]
	
	\centering
	\caption{Superficie sobre el paraboloide \( 3z = 9 - x^2 - y^2 \)}
	\begin{minipage}[t]{0.29\textwidth}
		\centering
		\includegraphics[width=\linewidth]{Img/41a.pdf}
	\end{minipage}
	\hspace{0.02\linewidth}
	\begin{minipage}[t]{0.29\textwidth}
		\centering
		\includegraphics[width=\linewidth]{Img/41b.pdf}
	\end{minipage}

	\fuentepropia
\end{figure}

Con la función ${\bf r}$ hallamos el siguiente vector.
\vspace{-5pt}
{\small
\begin{multline*}
\hspace*{-\multlinegap}
{\bf r}_{t} = \langle \sin(t)(2(1-u)\cos(t)-1-2u),\,
\cos(t)(2(u-1)\cos(t)+1+2u)-u+1, \\
\frac{2}{3}(u-1)\sin(t)(2u+(u-1)\cos(t)+1) \rangle \\[1ex]
{\bf r}_{u} = \langle (\cos(t)+2)\cos(t),\,
(\cos(t)+2)\sin(t),\,
\frac{2}{3}(\cos(t)+2)((1-u)\cos(t)-1-2u) \rangle
\end{multline*}
}

{\small
\begin{multline*}
{\bf r}_{t}\times{\bf r}_{u}=\bigg \langle-\frac{2}{3}\cos(t)(\cos(t)+2)(2u+(u-1)\cos(t)+1)^{2},\\
-\frac{2}{3}\sin(t)(\cos(t)+2)(2u+(u-1)\cos(t)+1)^{2},\\ -(\cos(t)+2)(2u+(u-1)\cos(t)+1)\bigg \rangle
\end{multline*}
}

De modo que el elemento diferencial de área está dado por:


{\small
\[
\lVert \mathbf{r}_t \times \mathbf{r}_u \rVert = \frac{1}{3}(\cos(t) + 2)(2u + (u - 1)\cos(t) + 1)f(u,t),
\]
}

en donde:
\begin{center}
	\resizebox{1.0\textwidth}{!}{%
		$f(u,t) = \sqrt{4\cos(t)(u - 1)((u - 1)\cos(t) + 4u + 2) + 13 + 16u(1 + u)}$%
	}
\end{center}


con lo cual, el área de la superficie es


\[\int_{0}^{1}\int_{0}^{2\pi} \|{\bf r}_{t}\times {\bf r}_{u}\|\,\, dtdu\approx 41.9822.\]

Para hacer los cálculos en \verb"Mathematica" definimos $x$ e $y$, que corresponden a las expresiones para la cardioide. La región entre las curvas se parametriza como \verb"r", mientras que
\verb"p" es la porción sobre el paraboloide cuya base está dada por \verb"r" y \verb"q", que definen la curva de intersección. El gráfico de la curva de intersección y la superficie se logran con las siguientes instrucciones:



\begin{tcolorbox}[breakable,enhanced,title=\bf Instrucción en \verb"Mathematica", top=2pt,bottom=2pt,boxsep=0pt,before skip=2pt,after skip=0pt]
\begin{Verbatim}[formatcom=\letraverbatim]
{x,y}={(1-Cos[t])*Cos[t],(1-Cos[t])*Sin[t]};
z=(9-x^2-y^2)/3;
r=(1-u)*{x,y}+u*{3*Cos[t],3*Sin[t]};
p={r[[1]],r[[2]],(9-r[[1]]^2-r[[2]]^2)/3};
q={x,y,z};


El gráfico de la curva de intersección y la superficie
se logran con las siguientes instrucciones


a=ParametricPlot3D[q,{t,0,2*Pi},PlotStyle->
{Black,Thickness->0.01},PlotPoints->80];
b=ParametricPlot3D[p,{t,0,2*Pi},{u,0,1},Mesh->False,
PlotStyle->{Orange,PlotPoints->120}];


Show[b,a]


Hallamos $\|{\bf r}_{t}\times {\bf r}_{u}\|$ 
que se denomina $A$ y se integra numéricamente,
A=Simplify[Norm[Cross[D[p,t],D[p,u]]]];
NIntegrate[A,{t,0,2*Pi},{u,0,1}]
\end{Verbatim}
\end{tcolorbox}




\subsection*{Un cilindro y dos planos}
Las ecuaciones $4x^{2}+z^{2}=4$, $z=3+2y$ y $z=5-x-3y$ representan un cilindro elíptico y dos planos, respectivamente.
Estas superficies acotan un sólido como se muestra a continuación:




\begin{figure}[H]
\centering
\caption{cilindro $4x^{2}+z^{2}=4$ acotado por dos planos}
\includegraphics[width=0.4\textwidth]{Img/55d.pdf} \hspace{5mm}
\includegraphics[width=0.4\textwidth]{Img/55c.pdf}

\includegraphics[width=0.4\textwidth]{Img/55b.pdf}
\fuentepropia
\end{figure}

Sobre este sólido parametrizado, las tres superficies que lo determinan y calcularemos el área de cada una.

La proyección del cilindro en el plano $xz$ corresponde a la ecuación $4x^{2}+z^{2}=4,$
que puede parametrizarse como $x=\cos(t),$ $z=2\sin (t),$ para $t\in \lbrack 0,2\pi ].$
La coordenada $y$ está determinada por las siguientes expresiones:

En el plano $z=3+2y$ obtenemos $y=\sin (t)-\frac{3}{2}$ y en el plano
$z=5-x-3y$ se sigue que $y=-\frac{2}{3}\sin (t)+\frac{5}{3}-\frac{1}{3}\cos (t).$

Definimos ahora las funciones vectoriales que determinan la intersección del cilindro con cada uno de los dos planos.
\begin{eqnarray*}
\overrightarrow{a}(t)&=&\la \cos(t),\sin(t)-\frac{3}{2},2\sin (t)\ra, \, t\in [0,2\pi] \\
\overrightarrow{b}(t)&=& \la \cos(t) ,-\frac{2}{3}\sin (t)+\frac{5}{3}-\frac{1}{3}\cos (t),2\sin (t)\ra , \, t\in [0,2\pi]
\end{eqnarray*}

\subsubsection*{Área sobre la superficie del cilindro}
Una combinación convexa de las funciones
$\overrightarrow{a}(t)$ y $\overrightarrow{b}(t)$ parametriza la superficie del cilindro que se encuentra entre los dos planos. Es decir,

\begin{multline*}
{\bf r}(u,t) = ( 1-u) \overrightarrow{a}(t)+u\overrightarrow{b}(t) \\
=\la \cos (t),\frac{1}{3}( 3-5u)\sin (t) +\frac{19}{6}u-\frac{1}{3}u\cos (t)-\frac{3}{2},2\sin (t)\ra,
\end{multline*}
con $t\in [0,2\pi],\,u\in [0,1].$
De aquí se obtienen las siguientes funciones.
\begin{eqnarray*}
{\bf r}_{t}&=& \la -\sin(t),\frac{1}{3}(3-5u)\cos(t)+\frac{1}{3}u\sin(t),2\cos(t)\ra \\
{\bf r}_{u}&=& \la 0,-\frac{5}{3}\sin(t)+\frac{19}{6}-\frac{1}{3}\cos(t),0\ra
\end{eqnarray*}

También se calcula
\begin{center}
\resizebox{1.0\textwidth}{!}{%
$\mathbf{r}_{t}\times\mathbf{r}_{u}= \left \langle2\cos(t)\left (\frac{5}{3}\sin(t)-\frac{19}{6}+\frac{1}{3}\cos(t)\right),0,\sin(t)\left (\frac{5}{3}\sin(t)-\frac{19}{6}+\frac{1}{3}\cos(t)\right) \right \rangle,$
	}
\end{center}

cuya magnitud es
\[\|{\bf r}_{t}\times{\bf r}_{u}\|=\frac{\sqrt{5+3\cos(2t)}(8+\cos(t)-4\sin (t))}{3\sqrt{2}};\]
y con esta expresión se obtiene el área, la superficie en mención.
\[\bigintss_{0}^{1}\bigintss_{0}^{2\pi}\frac{\sqrt{5+3\cos(2t)}(8+\cos(t)-4\sin(t))}{3\sqrt{2}} \, dtdu \approx 25.836.\]
Los cálculos anteriores se pueden realizar en \verb"Mathematica"
mediante la ejecución de las siguientes instrucciones:



\begin{tcolorbox}[breakable,enhanced,title=\bf Instrucción en \verb"Mathematica", top=2pt,bottom=2pt,boxsep=0pt,before skip=2pt,after skip=0pt]
\begin{Verbatim}[formatcom=\letraverbatim]
a={Cos[t],(-3+2*Sin[t])/3,2*Sin[t]};
b={Cos[t],(5-Cos[t]-2*Sin[t])/3,2*Sin[t]};
B=Norm[Cross[D[(1-u)*a+u*b,t],D[(1-u)*a+u*b,u]]]
A=Assuming[t>=0&&u>=0,Simplify[B]]
NIntegrate[A,{t,0,2Pi},{u,0,1}]
\end{Verbatim}
\end{tcolorbox}



\subsubsection*{Área sobre los planos}
La región interior a la elipse se parametriza como
$x=u\cos(t),$ $z=2u\sin (t),$ para $t\in \lbrack 0,2\pi ],u\in \lbrack 0,1].$
La coordenada $y$ se puede determinar reemplazando $x$ y $z$ en la ecuación de los planos, de esta manera para  $z=3+2y$ se sigue $y=u\sin (t)-\frac{3}{2},$
mientras que en $z=5-x-3y$ se obtiene que
$y=-\frac{2}{3}u\sin (t)+\frac{5}{3}-\frac{1}{3}u\cos (t).$

Por lo tanto, es natural que se definan las siguientes funciones que parame- \\ trizan la región sobre cada uno de los planos que acotan el cilindro.
\begin{eqnarray*}
\overrightarrow{\alpha }(t,u)&=&\la u\cos (t),u\sin (t)-\frac{3}{2},2u\sin (t)\ra \\
\overrightarrow{\beta }(t,u)&=&\la u\cos (t),-\frac{2}{3}u\sin (t)+\frac{5}{3}-\frac{1}{3}u\cos (t),2u\sin (t)\ra.
\end{eqnarray*}

El área sobre el plano $z=3+2y$ dentro del cilindro puede determinarse con el siguiente procedimiento: primero, se derivan parcialmente las funciones, según el caso.
\begin{eqnarray*}
\overrightarrow{\alpha }{t}&=& \la -u\sin (t),u\cos (t),2u\cos (t)\ra  \\
\overrightarrow{\alpha }{u}&=& \la \cos (t),\sin (t),2\sin (t)\ra.
\end{eqnarray*}
Luego se calcula el elemento diferencial de área 


\begin{align*}
	\|\mathbf{r}_t \times \mathbf{r}_u\| &= \| \langle -u\sin(t), u\cos(t), 2u\cos(t) \rangle \times \langle \cos(t), \sin(t), 2\sin(t) \rangle \| \\
	&= \sqrt{5}u
\end{align*}

Y, por último, este término se integra con los límites establecidos
\[\int_{0}^{2\pi }\int_{0}^{1}\| r_{t}\times r_{u}\|
dudt=\int_{0}^{2\pi }\int_{0}^{1}\sqrt{5}ududt= \sqrt{5}\pi.\]

Para la región sobre el plano con ecuación $z=5-x-3y,$
hallamos los mismos vectores que en el caso anterior
\begin{eqnarray*}
\overrightarrow{\beta}_{t}&=&  \la -u\sin (t),-\frac{1}{3}u( 2\cos (t)-\sin (t))  ,2u\cos (t) \ra  \\
\overrightarrow{\beta}_{u}&=& \la \cos (t),-\frac{2}{3}\sin (t)-\frac{1}{3}\cos (t),2\sin (t)\ra.
\end{eqnarray*}
El elemento diferencial de área en este caso está dado por
\[\| \overrightarrow{\beta}_{t} \times \overrightarrow{\beta}_{u} \|= \frac{2}{3}\sqrt{11}u\]
% $\left \| \la -u\sin (t),-\frac{1}{3}u( 2\cos (t)-\sin (t))  ,2u\cos (t)\ra  \times \la \cos (t),-\frac{2}{3}\sin (t)-\frac{1}{3}\cos (t),2\sin (t)\ra \right \| =  \frac{2}{3}\sqrt{11}u$
y el área de la superficie señalada es
\[\int_{0}^{2\pi }\int_{0}^{1}\| r_{t}\times r_{u}\|
dudt=\int_{0}^{2\pi }\int_{0}^{1}\frac{2}{3}\sqrt{11}ududt= \frac{2}{3}\pi \sqrt{11}.\]

\subsection*{Una porción de paraboloide}
El paraboloide con ecuación $\frac{(z-2)^2}{4}+\frac{(y-3)^2}{9}=1-\frac{x}{4}$
posee vértice en el punto $(4,3,2)$
y abre hacia la parte negativa del eje $x$; consideraremos la porción donde $x \geq 0$ y presentaremos además la proyección en el plano $yz$, que es una elipse.

Para rellenar el interior de la elipse se pueden usar las expresiones \linebreak
$y=3+3u\sin t,z=2+2u\cos t,$ donde $ u\in \lbrack 0,1]$ y $t \in \lbrack 0,2\pi ].$
La coordenada $x$ sobre el paraboloide se obtiene al sustituir $z$ y $y$ en dicha ecuación, es decir,

\begin{figure}[H]
	\centering
	\caption{Paraboloide y su proyección en el plano \( yz \)}
	\begin{minipage}[t]{0.23\textwidth}
		\centering
		
		\includegraphics[width=\linewidth]{Img/32.pdf}
	\end{minipage}
	\hspace{0.02\linewidth}
	\begin{minipage}[t]{0.27\textwidth}
		\centering
		
		\includegraphics[width=\linewidth]{Fig/131.pdf}
	\end{minipage}
	\hspace{0.02\linewidth}
	\begin{minipage}[t]{0.27\textwidth}
		\centering
		
		\includegraphics[width=\linewidth]{Fig/132.pdf}
	\end{minipage}
	\fuentepropia
\end{figure}

\[x=  4\left( 1-\frac{\left( 2u\cos (t) \right) ^{2}}{4}- \frac{\left( 3u\sin (t) \right) ^{2}}{9}\right) =4-4u^{2}.\]

Entonces, la función ${\bf r}(u,t)=\left\langle 4-4u^{2},3+3u\sin t,2+2u\cos
t\right\rangle ,$ $u\in \lbrack 0,1],\,\,t\in \lbrack 0,2\pi ], $
parametriza la superficie del paraboloide, con $x\geq 0.$
Con esta función hallamos los siguientes vectores
\begin{eqnarray*}
{\bf r}_{u} &=& \left\langle -8u,3\sin (t) ,2\cos (t) \right\rangle\\
{\bf r}_{t} &=& \left\langle 0,3u\cos (t) ,-2u\sin (t) \right\rangle.
\end{eqnarray*}
Entonces,
${\bf r}_{t}\times{\bf r}_{u}=\left\langle 6u,16u^{2}\sin (t),24u^{2}\cos (t) \right\rangle,$
por tanto,


\[\left\| {\bf r}_{u}\times {\bf r}_{t}\right\| = 2u\sqrt{9+64u^{2}+80u^{2}\cos ^{2}(t) }.\]
De esta manera, el área de la superficie del paraboloide está dada por
%$\int_{0}^{1}\int_{0}^{2\pi }\left\| {\bf r}_{u}\times {\bf r}_{t}\right\| dtdu  \approx 46.979.$
{\footnotesize
	\begin{equation*}
		\int_{0}^{1}\int_{0}^{2\pi} \|\mathbf{r}_{u} \times \mathbf{r}_{t}\| \, dt\,du =
		\int_{0}^{1}\int_{0}^{2\pi}
		2u\sqrt{9 + 64u^2 + 80u^2 \cos^2(t)} \, dt\,du \approx 46{,}979
	\end{equation*}
}
Este cálculo se hace en \verb"Mathematica" con las siguientes instrucciones



\begin{tcolorbox}[breakable,enhanced,title=\bf Instrucción en \verb"Mathematica", top=2pt,bottom=2pt,boxsep=0pt,before skip=2pt,after skip=0pt]
\begin{Verbatim}[formatcom=\letraverbatim]
y=3+3*u*Sin[t],
z=2+2*u*Cos[t]};
x=Simplify[4*(1-(z-2)^2/4-(y-3)^2/9)]
r={x,y,z};
A=Norm[Cross[D[r,t],D[r,u]]];
NIntegrate[A,{t,0,2Pi},{u,0,1}]
\end{Verbatim}
\end{tcolorbox}



\subsection*{Un cilindro cardióidico y un plano}
La superficie del cilindro cardiogénico con ecuación
$r=2(1+\cos(t))$, acotada por los planos $z=0$ y $z=2-x/2$, se presenta a continuación. También se muestra la porción sobre el plano $z=2-x/2$ acotada por el cilindro cardiogénico. Hallaremos el área de estas dos superficies:


\begin{figure}[H]
\centering
\caption{Superficies sobre el cilindro cardiogénico y sobre el plano \( z = 2 - \frac{x}{2} \)}
\includegraphics[width=0.4\textwidth]{Fig/134.pdf} \hspace{5mm}
\includegraphics[width=0.4\textwidth]{Fig/135.pdf}
\fuentepropia
\end{figure}

Sobre el plano $xy,$ la cardioide $r=2(1+\cos(t))$ se parametriza como
\[\la 2\left( 1+\cos (t) \right) \cos (t) ,2\left(1+\cos (t) \right) \sin (t),0 \ra.\]
En el plano con ecuación $z=2-x/2,$ se satisface que
\[z=2-2\left( 1+\cos (t) \right) \cos(t) /2= 2-\left( 1+\cos t\right) \cos t,\]
por tanto, la curva de intersección entre el plano y el cilindro cardiogénico se parametriza como
\[\la 2\left(1+\cos(t)\right)\cos(t),2\left(1+\cos(t)\right)\sin(t),2-(1+\cos (t))\cos(t) \ra,\]
con lo que forma una combinación convexa entre dos puntos cuyos primer y segundo componente son iguales, uno sobre el plano $z=0$ y el otro sobre el plano
$z=2-x/2.$ Se parametriza la superficie lateral, esto es
\begin{multline*}
{\bf r}(u,t)=\left( 1-u\right) \la 2\left( 1+\cos (t) \right) \cos \left(
t\right) ,2\left( 1+\cos (t) \right) \sin (t) ,0\ra \\
+u\la 2\left( 1+\cos (t) \right) \cos(t) ,2\left( 1+\cos (t) \right) \sin (t)
,2-\left( 1+\cos (t) \right) \cos (t) \\
= \la 2( \cos (t)+1) \cos (t),2\sin (t)+\sin (2t),u(2-\left( 1+\cos (t) \right) \cos (t)) \ra
\end{multline*}
%
$t\in [0,2\pi],$ $u\in [0,1].$ 
\\

Con esta función obtenemos los siguientes vectores:
\begin{eqnarray*}
{\bf r}_{t} &=& \big \langle-2\sin (t) \left( 2\cos (t) +1\right) ,2(\cos(t)+\cos ^{2}(t)-\sin^{2}(t)) , \\
&& \hskip 6cm u\sin (t) \left( 2\cos (t)+1\right) \big \rangle \\
{\bf r}_{u} &=& \la 0,0,\left( \cos (t) +2\right) \left( 1-\cos (t) \right) \ra\\
{\bf r}_{t}\times {\bf r}_{u}&=&\big\langle 2\sin^{2}(t)\left( 2\cos(t) -1\right)\left(\cos(t)+2\right) ,\\
&&\hskip1.5cm 2\sin(t) \left( 2\cos (t) +1\right) \left( \cos (t) +2\right) \left( 1-\cos (t) \right) ,0\big \rangle.
\end{eqnarray*}
Con lo cual,
\[\left\| {\bf r}_{t}\times {\bf r}_{u}\right\| = 2\sin (t) \left( \cos (t) +2\right) \sqrt{2-2\cos(t);}\]
entonces el área lateral del cilindro acotada por los dos planos está dada por

{\small
\[\int_{0}^{1}\int_{0}^{2\pi}\left|2\sin(t)\left(\cos(t)+2\right)\sqrt{2-2\cos(t)}\right|dtdu= \frac{96}{5}.\] 
}
El valor anterior se obtiene en \verb"Mathematica" con ayuda de las siguientes instrucciones:



\begin{tcolorbox}[breakable,enhanced,title=\bf Instrucción en \verb"Mathematica", top=2pt,bottom=2pt,boxsep=0pt,before skip=2pt,after skip=0pt]
\begin{Verbatim}[formatcom=\letraverbatim]
x=2*(1+Cos[t])*Cos[t);
y=2*(1+Cos[t]*Sin[t];
z=2-x/2;
r=(1-u)*{x,y,0}+u*{x,y,z);
A=Norm[Cross[D[r,t],D[r,u]]];
Integrate[A,{u,0,1},{t,0,2Pi}]
\end{Verbatim}
\end{tcolorbox}




Hallaremos ahora el área de la región sobre el plano $z=2-x/2,$ acotada por el cilindro cardiogénico.
El interior de la cardioide en el plano $xy$ se parametriza como
\[\la 2u\left( 1+\cos(t)\right) \cos(t) ,2u\left( 1+\cos(t)\right) \sin(t),0\ra,\]
con $t\in \lbrack 0,2\pi ],$ $u\in \lbrack 0,1].$
Sobre el plano $z=2-x/2,$ se cumple que
\[z=2-2u\left( 1+\cos(t)\right)\cos(t)/2=2-u\left(\cos t+1\right)\cos(t).\]

Entonces, la parte del plano acotada por el cilindro cardiogénico se parame- \\ triza por medio de la siguiente función
\[
\adjustbox{max width=\textwidth}{
	${\bf r}(u,t)=\big\langle 2u\left( 1+\cos(t)\right) \cos(t),2u\left(1+\cos(t)\right)\sin(t),2-u(\cos(t)+1)\cos(t) \big\rangle,$
}
\] con $t\in \lbrack 0,2\pi ],$ $u\in \lbrack 0,1].$ De aquí se obtienen los siguientes vectores:
\begin{align*}
	\mathbf{r}_t &= \left\langle -2u\sin(t)(2\cos(t) + 1), 2u\left(\cos(t) - \sin^2(t) + \cos^2(t)\right), \right. \\
	&\quad \left. u\sin(t)(2\cos(t) + 1) \right\rangle \\
	\mathbf{r}_u &= \left\langle 2(1 + \cos(t))\cos(t), 2(1 + \cos(t))\sin(t), -(1 + \cos(t))\cos(t) \right\rangle \\
	\mathbf{r}_t \times \mathbf{r}_u &= \left\langle -2u(1 + \cos(t))^2, 0, -4u(1 + \cos(t))^2 \right\rangle \\
	&= -2u(1 + \cos(t))^2 \langle 1, 0, 2 \rangle
\end{align*}
Estos vectores, a su vez, nos permiten obtener el elemento diferencial de área:
$\left\|  {\bf r}_{t}\times {\bf r}_{u}\right\|= 2\sqrt{5}\left|u\left( 1+\cos(t)\right)^{2}\right| .$ Por lo tanto, el área de la superficie en mención está dada por
\[\int_{0}^{1}\int_{0}^{2\pi }2\sqrt{5}\left| u\left( 1+\cos (t) \right) ^{2}\right|
dtdu= 3\pi \sqrt{5}.\]

En \verb"Mathematica" este valor se consigue con las siguientes instrucciones.



\begin{tcolorbox}[breakable,enhanced,title=\bf Instrucción en \verb"Mathematica", top=2pt,bottom=2pt,boxsep=0pt,before skip=2pt,after skip=0pt]
\begin{Verbatim}[formatcom=\letraverbatim]
x=2*u*(1+Cos[t])*Cos[t;
y=2*u*(1+Cos[t])*Sin[t];
z=2-x/2;
r=(1-u)*{x,y,0}+u*{x,y,z};
A=Simplify[Norm[Cross[D[r,t],D[r,u]]]];
NIntegrate[A,{u,0,1},{t,0,2Pi}]
\end{Verbatim}
\end{tcolorbox}


%\clearpage

\subsection*{Dos cilindros cardiotónicos entre dos planos}
Las ecuaciones $r = 1-\cos(t)$ y $r = 2- 2\cos(t)$ representan, en el plano, dos cardioides, y en el espacio estas mismas ecuaciones representan dos cilindros cardiotónicos \footnote{Este término fue acuñado por el profesor Juan Zambrano de la Universidad Distrital Francisco José de Caldas.}. Estas cilíndricas se cortan por los planos $z=0$ y $z=2-x/2$.
%
\begin{figure}[H]
	\centering
	\caption{Superficie sobre el plano \( z = 2 - x/2 \) acotada por dos cilindros cardióidicos}
	\begin{minipage}[t]{0.28\linewidth}
		\centering
		\includegraphics[width=\linewidth]{Fig/136.pdf}
	\end{minipage}
	\hspace{0.5em}
	\begin{minipage}[t]{0.28\linewidth}
		\centering
		\includegraphics[width=\linewidth]{Img/36b.pdf}
	\end{minipage}
	\vspace{1ex}
	\fuentepropia
\end{figure}
%
En el plano $xy$ los cilindros cardiotónicos se  parametrizan en la forma:
\begin{eqnarray*}
\overrightarrow{\alpha} (t) &=& \la(1-\cos(t))\cos(t), (1-\cos(t))\sin(t), 0\ra \\
\overrightarrow{\beta} (t)& =& \la(2-2\cos(t))\cos(t), (2-2\cos(t))\sin(t), 0\ra  ,\,\, t \in[0, 2\pi]
\end{eqnarray*}
Ahora formamos una combinación convexa entre estas, que determina la región contenida entre las dos curvas. Esto es, para $u\in[0, 1]$ y $t \in[0, 2\pi]$,
cada una de las curvas $\overrightarrow{\alpha }(t)$ y $\overrightarrow{\beta }(t)$
se proyecta sobre el plano $z=2-x/2$, de la siguiente forma:
\[
\adjustbox{max width=\textwidth}{
$\begin{array}{rcl}
\overrightarrow{p}(t) &=& \big\langle \left( 1-\cos(t) \right)\cos(t),\left( 1-\cos(t) \right)\sin(t), 2-\left( \left( 1-\cos(t) \right)\cos(t) \right)/2 \big\rangle \\[1ex]
\overrightarrow{q}(t) &=& \big\langle \left( 2-2\cos(t) \right)\cos(t), \left( 2-2\cos(t) \right)\sin(t), 2-\left( 2\left( 1-\cos(t) \right)\cos(t) \right)/2 \big\rangle
\end{array}$}
\]

\begin{figure}[H]
	\centering
	\caption{Cilindros cardiotónicos acotados por dos planos}
	\begin{minipage}[t]{0.28\linewidth}
		\centering
		
		\includegraphics[width=\linewidth]{Fig/137.pdf}
	\end{minipage}
	\hspace{1em}
	\begin{minipage}[t]{0.28\linewidth}
		\centering
		
		\includegraphics[width=\linewidth]{Fig/138.pdf}
	\end{minipage}
	\vspace{1ex}
	\fuentepropia
\end{figure}

Con el propósito de parametrizar la parte sobre los cilindros entre los planos $z=0$ y $z=2-x/2$, se hará una combinación convexa entre las expresiones $\overrightarrow{p}(t)$ y $\overrightarrow{\alpha}(t)$ por un lado
y $\overrightarrow{q}(t)$ con $\overrightarrow{\beta}(t)$ por el otro.




\begin{figure}[H]
\centering
\caption{
	Cilindros cardiotónicos
}
\includegraphics[width=0.4\textwidth]{Img/57a.pdf} \hspace{5mm}
\includegraphics[width=0.4\textwidth]{Img/57b.pdf}
\fuentepropia
\end{figure}

Estas figuras se obtienen en \verb"Mathematica" con las siguientes instrucciones

 % espacio entre texto y código



\begin{tcolorbox}[breakable,enhanced,title=\bf Instrucción en \verb"Mathematica", top=2pt,bottom=2pt,boxsep=0pt,before skip=2pt,after skip=0pt]
\begin{Verbatim}[formatcom=\letraverbatim]
ParametricPlot3D[{{2*Cos[t]*(1-Cos[t]),
2*Sin[t]*(1-Cos[t]),u(2-Cos[t]+Cos[t]^2)},
{Cos[t]*(1-Cos[t]),Sin[t]*(1-Cos[t]),
(u/2)*(4-Cos[t]+Cos[t]^2)}},{u,0,1},{t,0,2*Pi},
PlotStyle->{Orange,RGBColor[0.3,0.8,0.4]},
Ticks->{{-2,0,1},{-2,0,1},{-1,0,1}},
Mesh->False,PlotPoints->80]
\end{Verbatim}
\end{tcolorbox}




Ahora calcularemos el área sobre los cilindros.

\subsubsection*{Cilindro menor}
Para $\overrightarrow{p}(t)$ y $\overrightarrow{\alpha }(t)$ se puede simplificar la expresión
$(1-u)\overrightarrow{\alpha }(t)+u\overrightarrow{p}(t),$
la cual proporciona la función que describe la parte del cilindro menor y que se halla entre los planos $z=0$ y $z=2-x/2$.


\begin{multline*}
{\bf r}\left( u,t\right) = \left( 1-u\right) \left\langle \left( 1-\cos (t) \right) \cos
(t) ,\left( 1-\cos (t) \right) \sin (t),0\right\rangle \\
+u\left\langle \left( 1-\cos (t) \right) \cos (t) ,\left( 1-\cos (t) \right) \sin (t).
2-\left( \left( 1-\cos (t) \right) \cos (t) \right)
/2\right\rangle \\
=\left\langle \cos (t) \left( 1-\cos (t) \right)
\sin (t) \left( 1-\cos (t) \right) ,\frac{1}{2}u\left( 4-\cos (t) +\cos ^{2}(t) \right)
\right\rangle,
\end{multline*}
en donde $t\in \lbrack 0,2\pi ],$ $u\in \lbrack 0,1].$
Con el fin de hallar el área superficial, se requieren los siguientes elementos.


\begin{multline*}
{\bf r}_{u}= \la 0,0,2-\frac{1}{2}\cos (t) +\frac{1}{2}\cos ^{2}(t) \ra \\
{\bf r}_{t}= \big\langle \sin (t) \left( 2\cos (t) -1\right) ,\left( 2\cos (t) +1\right) \left( 1-\cos (t) \right) , \hskip1.5cm \\ \frac{1}{2}u\sin (t) \left( 1-2\cos (t) \right) \big\rangle.
\end{multline*}

El elemento diferencial de área es
\begin{equation}
\left\| \mathbf{r}_{t}\times \mathbf{r}_{u}\right\| = \frac{1}{2}\left( \cos ^{2}(t) -\cos (t) +4\right) \sqrt{2-2\cos (t)}.
\end{equation}
Entonces, el área superficial sobre el cilindro menor es
\begin{equation}
\int_{0}^{1}\int_{0}^{2\pi }\frac{1}{2}\left(\cos ^{2}(t)
-\cos (t) +4\right) \sqrt{2-2\cos (t)}dtdu= \frac{96}{5}.
\end{equation}
Estos cálculos se hacen muy simples, ejecutando las siguientes instrucciones:



\begin{tcolorbox}[breakable,enhanced,title=\bf Instrucción en \verb"Mathematica", top=2pt,bottom=2pt,boxsep=0pt,before skip=2pt,after skip=0pt]
\begin{Verbatim}[formatcom=\letraverbatim]
{x,y}={Cos[t]*(1-Cos[t]),Sin[t]*(1-Cos[t])};
a={x,y,0};
p={x,y,2-x/2};
r=(1-u)*a+u+p;
Integrate[Norm[Cross[D[r,u]]],{t,0,2*Pi},{u,0,1}]
\end{Verbatim}
\end{tcolorbox}



\subsubsection*{Cilindro mayor}
Para $\overrightarrow{q}(t)$ y $\overrightarrow{\beta}(t)$, es necesario simplificar la expresión 
\[
(1 - u)\overrightarrow{\beta}(t) + u\overrightarrow{q}(t),
\]
lo que nos provee de la función que parametriza la región entre los planos sobre el cilindro mayor, esto es

\begin{flushleft}
\adjustbox{max width=\textwidth}{
$\begin{array}{l}
\hspace{-0.6em}\mathbf{r}(u,t) = (1 - u)(2(1 - \cos(t))\cos(t),\; 2(1 - \cos(t))\sin(t),\; 0) \\[1ex]
\quad + u((2 - 2\cos(t))\cos(t),\; (2 - 2\cos(t))\sin(t),\; 2 - (2(1 - \cos(t))\cos(t))/2) \\[1ex]
= \langle 2\cos(t)(1 - \cos(t)),\; 2\sin(t)(1 - \cos(t)),\; u(2 - \cos(t) + \cos^2(t)) \rangle
\end{array}$}
\end{flushleft}

Con lo anterior, se halla el elemento diferencial de área:
\[
\left\| {\bf r}_{t} \times {\bf r}_{u} \right\| = \left| 2\left(2 - \cos(t) + \cos^2(t)\right) \sqrt{2 - 2\cos(t)} \right|.
\]

Luego, el área superficial del cilindro mayor está dada por:
\[
\int_{0}^{1} \int_{0}^{2\pi} \left| 2\left(2 - \cos(t) + \cos^2(t)\right) \sqrt{2 - 2\cos(t)} \right| \, dt \, du = \frac{224}{5}.
\]

En \verb"Mathematica", estos cálculos se hacen con las siguientes instrucciones.



\begin{tcolorbox}[breakable,enhanced,title=\bf Instrucción en \verb"Mathematica", top=2pt,bottom=2pt,boxsep=0pt,before skip=2pt,after skip=0pt]
\begin{Verbatim}[formatcom=\letraverbatim]
{x,y}={2*Cos[t]*(1-Cos[t]),2*Sin[t]*(1-Cos[t])};
r=(1-u)*{x,y,0}+u*{x,y,2-x/2};
Integrate[Norm[Cross[D[r,t,[r,u]]],{t,0,2*Pi},{u,0,1}]
\end{Verbatim}
\end{tcolorbox}



\subsection*{Cono oblicuo truncado}
Consideremos las funciones vectoriales $\overrightarrow{\alpha }(t) =\langle 1+\cos \left(
t\right) ,\sin (t) ,0\rangle$ y $\overrightarrow{\beta }(t) =\langle 2+2\cos \left(
t\right) ,2\sin (t) ,3 \rangle,$
que representan dos circunferencias en el espacio. Una combinación convexa de estas define la función r:

\begin{flalign*}
	&\mathbf{r}(t, u) = (1-u)\overrightarrow{\alpha}(t) + u\overrightarrow{\beta}(t) = \small{\left\langle (1+\cos(t))(1+u), \sin(t)(1+u), 3u\right\rangle}, &&
\end{flalign*}

con $t\in \lbrack 0,2\pi ],$ $u\in \lbrack 0,1].$ Esta función parametriza el cono truncado oblicuo que se muestra en la figura.



\begin{figure}[H]
\centering
\caption{Cono oblicuo truncado}
\includegraphics[width=0.5\textwidth]{Img/19a.pdf}
\fuentepropia
\end{figure}

Con la función ${\bf r}\left( t,u\right) $ hallamos${\bf r}_{t}=\left\langle -\sin (t) \left( 1+u\right) ,\cos \left(
t\right) \left( 1+u\right) ,0\right\rangle \,$ \break
y $\,{\bf r}_{u}=\left\langle 1+\cos (t) ,\sin (t),3\right\rangle,$
que determinan el elemento diferencial de área
{\small
\begin{align*}
	\|\mathbf{r}_t \times \mathbf{r}_u\| &= \left\| \left\langle 3\cos(t)(1+u), 3\sin(t)(1+u), -(1+\cos(t))(1+u) \right\rangle \right\| \\
	&= (1+u)\sqrt{\cos^2(t) + 10 + 2\cos(t)},
\end{align*}
}
cuya integral proporciona el área de esta región
\[\int_{0}^{2\pi }\int_{0}^{1}\left( 1+u\right) \sqrt{\cos ^{2}(t) +10+2\cos (t) }dudt \approx 30.4679.\]
La figura y los cálculos se obtienen con las siguientes instrucciones:



\begin{tcolorbox}[breakable,enhanced,title=\bf Instrucción en \verb"Mathematica", top=2pt,bottom=2pt,boxsep=0pt,before skip=2pt,after skip=0pt]
\begin{Verbatim}[formatcom=\letraverbatim]
ParametricPlot3D[{(1+u)*(1+Cos[t]),(1+u)*Sin[t],3*u},
{u,0,1},{t,0,2*Pi},PlotPoints->80,Mesh->True,
PlotStyle->Blend[{Red,Black,Yellow}]]
{a,b}={{1+Cos[t],Sin[t],0},{2+2*Cos[t],2*Sin[t],3}}
A=Norm[Cross[D[(1-u)*a+u*b,t],D[(1-u)*a+u*b,u]]]
NIntegrate[A,{t,0,2*Pi},{u,0,1}]
\end{Verbatim}
\end{tcolorbox}



\subsection*{Superficie obtenida entre dos elipses}

Las funciones vectoriales
$\overrightarrow{a}(t)= \la 2\cos(t),\sin(t),0\ra$ y $\overrightarrow{b}(t)= \big \langle \cos(t), \linebreak 3\sin(t),3\big\rangle$,
con $t \in [0,2\pi]$ representan dos elipses en el espacio.
Con la función
\begin{eqnarray*}
{\bf r}\left( t,u\right) &=&\left( 1-u\right) \left\langle 2\cos (t),\sin (t) ,0\right\rangle +u\left\langle \cos (t),3\sin (t) ,3 \right\rangle \\
&=&  \left\langle \cos (t)(2-u),\sin (t)(1+2u),3u\right\rangle,
\end{eqnarray*}
Con $u \in [0,1]$ y $t \in [0,2\pi]$, se parametriza la superficie que se muestra en la siguiente figura.

%


\begin{figure}[H]
\centering
\caption{Superficie que conecta dos elipses.}
\includegraphics[width=0.4\textwidth]{Fig/139.pdf} \hspace{5mm}
\includegraphics[width=0.4\textwidth]{Img/20a.pdf}

\includegraphics[width=0.4\textwidth]{Img/20b.pdf}
\fuentepropia
\end{figure}

Este gráfico se obtiene con las siguientes instrucciones:



\begin{tcolorbox}[breakable,enhanced,title=\bf Instrucción en \verb"Mathematica", top=2pt,bottom=2pt,boxsep=0pt,before skip=2pt,after skip=0pt]
\begin{Verbatim}[formatcom=\letraverbatim]
ParametricPlot3D[{(2-u)*Cos[t],(1+2*u)*Sin[t],3*u},
{u,0,1},{t,0,2*Pi},PlotPoints->80,Mesh->False]
\end{Verbatim}
\end{tcolorbox}



Con ayuda de la función $\mathbf{r}$ se obtienen los siguientes vectores:\\
$\mathbf{r}_{u} =\left\langle -\cos(t),\, 2\sin(t),\, 3 \right\rangle$ y $\mathbf{r}_{t} = \left\langle -\sin(t)(2 - u),\, \cos(t)(1 + 2u),\, 0 \right\rangle$, los cuales describen las direcciones tangentes de la superficie parametrizada. Además, el producto cruzado\\ $\mathbf{r}_{t} \times \mathbf{r}_{u} = \left\langle 3\cos(t)(1 + 2u),\, 3\sin(t)(2 - u),\, 5\cos^2(t) + 2u - 4 \right\rangle$ nos proporciona el vector normal y define el elemento diferencial de área.
\[
\left\| {\\bf r}_{t} \times {\\bf r}_{u} \right\| = \sqrt{ \cos^2(t)\left(25\cos^2(t) - 67 + 92u + 27u^2\right) + 13(u - 2)^2 }.
\]

La integral que proporciona el área de la superficie en mención es:
\[
\int_{0}^{1} \int_{0}^{2\pi} \sqrt{ \cos^2(t)\left(25\cos^2(t) + P(u)\right) + 13(u - 2)^2 } \, dt \, du \approx 35.72,
\]

Estos cálculos se hacen en \verb"Mathematica" usando las siguientes sentencias.



\begin{tcolorbox}[breakable,enhanced,title=\bf Instrucción en \verb"Mathematica", top=2pt,bottom=2pt,boxsep=0pt,before skip=2pt,after skip=0pt]
\begin{Verbatim}[formatcom=\letraverbatim]
r=Simplify[(1-u)*{2*Cos[t],Sin[t],0}+u*{Cos[t],
3*Sin[t],3}];
A=Simplify[Norm[Cross[D[r,t],D[r,u]]]];
NIntegrate[A,{t,0,2Pi},{u,0,1}]
\end{Verbatim}
\end{tcolorbox}




\index{Cono cardioidico!oblicuo y truncado}
\subsection*{Cono cardiogénico oblicuo y truncado} 
A partir de las cardioides definidas por las funciones vectoriales
\begin{eqnarray*}
\overrightarrow{\alpha }(t) & \left\langle \left( 1+\cos \left(
t\right) \right) \cos (t) ,\left( 1+\cos (t) \right)
\sin (t) ,0\right\rangle \\
\overrightarrow{b}(t) & \left\langle 2\left( 1+\cos \left(
t\right) \right) \cos (t) ,2\left( 1+\cos (t)\right) \sin (t) ,3\right\rangle.
\end{eqnarray*}
Con $t \in [0,2\pi],$ se trazan segmentos rectilíneos entre estas dos curvas, lo que produce una parametrización de un cono truncado oblicuo, cuyos cortes paralelos al plano $xy$ son cardioides.
La función que parametriza la superficie obtenida y su gráfico son las siguientes:
\begin{eqnarray*}
{\bf r}\left( t,u\right) &=& (1-u)\overrightarrow{\alpha }(t)+ u \overrightarrow{\beta}(t) \\
&& \left\langle (1+\cos(t))(1+u) \cos (t),( 1+\cos (t))(1+u) \sin (t) ,3u  \right\rangle.
\end{eqnarray*}




\begin{figure}[H]
\centering
\caption{Cono cardiogénico oblicuo truncado}
\includegraphics[width=0.4\textwidth]{Fig/140.pdf} \hspace{5mm}
\includegraphics[width=0.4\textwidth]{Img/21a.pdf}

\includegraphics[width=0.4\textwidth]{Img/21b.pdf}
\fuentepropia
\end{figure}

La superficie se puede obtener ejecutando las siguientes instrucciones.



\begin{tcolorbox}[breakable,enhanced,title=\bf Instrucción en \verb"Mathematica", top=2pt,bottom=2pt,boxsep=0pt,before skip=2pt,after skip=0pt]
\begin{Verbatim}[formatcom=\letraverbatim]
ParametricPlot3D[{(2-u)*Cos[t]*(1+Cos[t]),
(2-u)*Sin[t]*(1+Cos[t]),3-3*u},{u,0,1},{t,0,2*Pi},
PlotPoints->90,Mesh->False]
\end{Verbatim}
\end{tcolorbox}



Con la función ${\bf r}$, se obtienen los siguientes vectores\\
${\bf r}_{t}= \left\langle -\sin (t) \left( 1+2\cos (t) \right) (1+u) ,\left( 1+\cos(t) \right)
\left( 2\cos(t) -1\right)( 1+u) ,0\right\rangle $\\
$ {\bf r}_{u}=\left\langle \cos (t) \left( 1+\cos \left(
t\right) \right) ,\sin (t) \left( 1+\cos (t) \right),3\right\rangle .$
Además,
\begin{multline*}
{\bf r}_{t}\times {\bf r}_{u} = (1+u) \big\langle 3\left( 1+\cos (t) \right) \left( 2\cos (t) -1\right),3\sin (t) \left( 1+2\cos (t) \right)  , \\
-\left( 1+\cos (t) \right) ^{2} \big \rangle
\end{multline*}

La norma de este último vector es el elemento diferencial de área
\[\left\| {\bf r}_{t}\times {\bf r}_{u}\right\| = (1+u) \sqrt{(1+\cos(t))(\cos ^{3}(t) +3\cos^{2}(t)+3\cos (t) +19)},\]
que, al integrarse, nos proporciona el área de la superficie que se considera
{\small
\[\int_{0}^{2\pi }\int_{0}^{1}\left\| {\bf r}_{t}\times {\bf r}_{u}\right\|  dudt \approx 39.4.\]
}

Este valor se obtiene ejecutando las siguientes instrucciones:



\begin{tcolorbox}[breakable,enhanced,title=\bf Instrucción en \verb"Mathematica", top=2pt,bottom=2pt,boxsep=0pt,before skip=2pt,after skip=0pt]
\begin{Verbatim}[formatcom=\letraverbatim]
r=Simplify[(1-u)*{(1+Cos[t])*Cos[t],
(1+Cos[t])*Sin[t],0}+u*{2*(1+Cos[t])*Cos[t],
2*(1+Cos[t])*Sin[t],3}];
A=Simplify[Norm[Cross[D[r,t],D[r,u]]]]
NIntegrate[A,{t,0,2Pi},{u,0,1}]
\end{Verbatim}
\end{tcolorbox}



\subsection*{Una flor sobre una esfera}
La ecuación $r=\sin(3\theta) $ representa, en coordenadas polares, una rosa de 3 pétalos; la proyectamos sobre la esfera centrada en el origen de radio $\sqrt{2}$.



\begin{figure}[H]
\centering
\caption{La flor $r=3\sin(3\theta)$ sobre una esfera}
\includegraphics[width=0.4\textwidth]{Fig/141.pdf} \hspace{5mm}
\includegraphics[width=0.4\textwidth]{Img/48a.pdf}
\fuentepropia
\end{figure}

En el plano, la flor (con su interior) se parametriza en  forma \\ {\small $x=u\sin (3t) \cos (t) ,$} {\small $y=u\sin (3t) \sin (t)$}
al reemplazar en la ecuación de la esfera,
{\small $x^{2}+y^{2}+z^{2}=2$,} se sigue que
{\small $z=\sqrt{2-u^{2}\sin^{2}(3t)},$}
entonces, la función
{\small ${\bf r}\left( t,u\right) =\left\langle u\sin (3t) \cos (t) ,u\sin (3t) \sin (t) ,\sqrt{2-u^{2}\sin^{2}(3t) }\right\rangle$}
con {\small $u \in [0,1]$, $t \in [0,2\pi]$,}
parametriza la región sobre la esfera cuya proyección en el plano $xy$ es la flor mencionada.
Los gráficos se obtienen ejecutando las siguientes instrucciones:



\begin{tcolorbox}[breakable,enhanced,title=\bf Instrucción en \verb"Mathematica", top=2pt,bottom=2pt,boxsep=0pt,before skip=2pt,after skip=0pt]
\begin{Verbatim}[formatcom=\letraverbatim]
A=ContourPlot3D[x^2+y^2+z^2==2,{x,-1.5,1.5},
{y,-1.5,1.5},{z,-1.5,1.5},ContourStyle->{Orange,
Opacity[0.6]},Mesh->False];
r={u*Sin[3*t]*Cos[t],u*Sin[3*t]*Sin[t]};
p={r[[1]],r[[2]],Sqrt[2-r[[1]]^2-r[[2]]^2]};
b=ParametricPlot3D[p,{u,0,1},{t,0,2*Pi},
PlotRange->{{-1,1},{-1,1},{1,1.5}},Mesh->False,
PlotStyle->RGBColor[0.8,0.3,0.3],PlotPoints->90,
Ticks->{{-1,0,1},{-1,0,1},{1.2,1.4}}];
c=ParametricPlot3D[{Sin[3*t]*Cos[t],Sin[3*t]*Sin[t],
Sqrt[2-Sin[3*t]^2]},{t,0,2*Pi},PlotPoints->80,
PlotStyle->{Black,Thickness->0.012}];
Show[A,c,b]
\end{Verbatim}
\end{tcolorbox}





\begin{figure}[H]
\centering
\caption{La flor $r=3\sin(3\theta)$ sobre una esfera}
\includegraphics[width=0.4\textwidth]{Img/68a.pdf} \hspace{5mm}
\includegraphics[width=0.4\textwidth]{Img/68b.pdf}
\fuentepropia
\end{figure}

Con la función ${\bf r}$ definida antes, hallamos los siguientes vectores:
{\small
\begin{eqnarray*}
{\bf r}_{t} &=& u\bigg\langle \cos(2t)+2\cos(4t),2\sin(4t)-\sin(2t),\frac{-3u\sin(6t)}{2\sqrt{2-u^{2}\sin^{2}(3t)  }} \bigg\rangle \\
{\bf r}_{u} &=& \sin (3t)\left\langle  \cos (t),
\sin (t) ,-\frac{u\sin (3t) }{\sqrt{2-u^{2}\sin^{2}(3t) }}\right\rangle
\end{eqnarray*}
}

{\small
\begin{flalign*}
	&\mathbf{r}_t \times \mathbf{r}_u = \left\langle -\frac{u^2 \cos(t) \sin^3(3t)}{\sqrt{2 - u^2 \sin^2(3t)}}, -\frac{u^2 \sin(t) \sin^3(3t)}{\sqrt{2 - u^2 \sin^2(3t)}}, -u \sin^2(3t) \right\rangle. &
\end{flalign*}
}

Entonces,
{\small
\[\left\| {\bf r}_{t}\times {\bf r}_{u}\right\|=\frac{\sqrt{2}u\sin^{2}(3t)}{\sqrt{2-u^{2}\sin^{2}(3t)}},\]
}
cuya integral nos proporciona claridad sobre la región considerada.
{\small
\[\bigintss_{0}^{2\pi}\bigintss_{0}^{1}\frac{\sqrt{2}u\sin ^{2}(3t) }{\sqrt{2-u^{2}\sin^{2}(3t) } }dudt \approx 1. 761.\]
}
La integral anterior se evalúa numéricamente con las siguientes sentencias.




\begin{tcolorbox}[breakable,enhanced,title=\bf Instrucción en \verb"Mathematica", top=2pt,bottom=2pt,boxsep=0pt,before skip=2pt,after skip=0pt]
\begin{Verbatim}[formatcom=\letraverbatim]
r={u*Sin[3*t]*Cos[t],u*Sin[3*t]*Sin[t]};
p={r[[1]],r[[2]],Sqrt[2-r[[1]]^2-r[[2]]^2]};
A=Norm[Cross[D[p,t],D[p,u]]];
NIntegrate[A,{t,0,2*Pi},{u,0,1}]
\end{Verbatim}
\end{tcolorbox}



\subsection*{Toros cardioidicos} \index{Toro cardioidico}
En esta sección consideraremos la superficie que se obtiene al hacer girar una cardioide alrededor de una recta en el mismo plano que no interseca a la figura.




\begin{figure}[H]
\centering
\caption{Cardioides}
\includegraphics[width=0.4\textwidth]{Fig/143.pdf} \hspace{5mm}
\includegraphics[width=0.4\textwidth]{Fig/144.pdf}

\includegraphics[width=0.4\textwidth]{Fig/145.pdf}
\fuentepropia
\end{figure}

\subsubsection*{Toro cardióidico generado por \boldmath$r=1-\sin(t)$ \unboldmath}
En el plano, la figura de la ecuación $r=1-\sin(t)$ representa una cardioide.
Esta figura tiene área determinada por
\[\int_{0}^{2\pi }\int_{0}^{1-\sin \theta }rdrd\theta = \frac{3}{2}\pi\]
y su longitud de arco se puede hallar como sigue
\begin{eqnarray*}
\int_{0}^{2\pi }\sqrt{r^{2}+\left( \frac{dr}{d\theta }\right) ^{2}}d\theta &=&
\int_{0}^{2\pi }\sqrt{\left( 1-\sin \theta \right) ^{2}+\left( -\cos \left(
\theta \right) \right) ^{2}}d\theta\\
&=&\int_{0}^{2\pi }\sqrt{2-2\sin \left(\theta \right) }d\theta = 8.
\end{eqnarray*}

La figura de la ecuación $r=1-\sin(t)$ se puede parametrizar en el plano en la forma
$\langle (1-\sin(t))\cos(t),(1-\sin(t))\sin(t) \rangle $ con $t\in [0,2\pi].$
Ahora, le sumaremos $3$ a la coordenada $x$, obteniendo la función $\overrightarrow{\alpha}(t)$,
y su gráfico se hará girar dicha figura alrededor de la recta $\theta = \pi/2.$


\begin{figure}[H]
\centering
\caption{Cardioide $\boldsymbol\alpha(t)=\langle 3+(1-\sin(t))\cos(t),(1-\sin(t))\sin(t)\rangle$}
\includegraphics[width=0.8\textwidth]{Fig/146.pdf}
\fuentepropia
\end{figure}

El área de una superficie de revolución está dada por
\begin{equation}\label{toros1}
A=\int_{a}^{b}x(t)\sqrt{\left( \frac{dx}{dt}\right) ^{2}+\left( \frac{dy}{dt}\right) ^{2}}dt
\end{equation}
Si el giro se hace alrededor del eje $y$, en este caso se obtiene

\[2\pi \int_{0}^{2\pi }\left( 3+\left( 1-\sin (t) \right) \cos
(t) \right) \sqrt{2-2\sin (t) }dt= 48\pi.\]

Para hallar el centro de masa de la curva, determinada por la expresión
$r=1-\sin (t) ,$ nótese que el término \[\sqrt{\left( \frac{dx}{dt}\right) ^{2}+\left( \frac{dy}{dt}\right) ^{2}}=\sqrt{2-2\sin (t) }.\]
De esta manera, las coordenadas del centro de masa de la curva están dadas por
\begin{eqnarray*}
\overline{x}&=&\frac{1}{8}\int_{0}^{2\pi }\left( 3+\left( 1-\sin \left(
t\right) \right) \cos (t) \right) \sqrt{2-2\sin (t) }dt= 3 \\
\overline{y} &=& \frac{1}{8}\int_{0}^{2\pi }\left( 1-\sin (t)\right) \sin (t) \sqrt{2-2\sin (t) }dt= -\frac{4}{5}.
\end{eqnarray*}

Usando el teorema de Pappus (\ref{teoremadepappus}), ver página (\pageref{teoremadepappus}), el área de la superficie de revolución de la curva $r=1-\sin(t)$ alrededor de la recta  $\theta = \pi /2$ es.
$2\pi \left( 3\right) \left( 8\right) = 48\pi .$ \index{Teorema de Pappus!área superficial}

Al ubicar la curva de giro en el espacio tridimensional y hacer coincidir el eje de giro con el eje $z$,
se satisface que si un punto sobre esta curva gira alrededor del eje $z$, entonces se obtiene una circunferencia y se genera una superficie que denominaremos \emph{toro cardioidico.}
\index{Toro cardioidico}




\begin{figure}[H]
\centering
\caption{Toro cardiogénico construido a partir de $r=1-\sin(t)$}
\includegraphics[width=0.4\textwidth]{Fig/147.pdf} \hspace{5mm}
\includegraphics[width=0.4\textwidth]{Img/63a.pdf}
\fuentepropia
\end{figure}

Al girar estos puntos sobre el eje $z$, todos los puntos conservan la misma coordenada en $z$,
y la coordenada en $x$ e $y$ se le agregará una $\cos(u)$ y $\sin(u)$ respectivamente, siendo $u$ un parámetro entre $0$ y $2 \pi$. Por lo tanto, esta superficie se parametriza como
\begin{flushleft} % Fuerza todo el bloque a la izquierda.
	\resizebox{1.0\textwidth}{!}{%
		$\mathbf{r}(t,u) = \left\langle (3+(1-\sin(t))\cos(t))\cos(u),(1-\sin(t))\sin(t)\sin(u),(1-\sin(t))\sin(t) \right\rangle$
	}
\end{flushleft}
con $t \in [0,2\pi],$ $u \in [0,2\pi].$




\begin{figure}[H]
\centering
\caption{Toro cardiogénico construido a partir de $r=1-\sin(t)$}
\includegraphics[width=0.4\textwidth]{Img/63b.pdf} \hspace{5mm}
\includegraphics[width=0.4\textwidth]{Img/63c.pdf}
\fuentepropia
\end{figure}

Las figuras anteriores se obtienen en \verb"Mathematica" ejecutando las siguientes instrucciones:



\begin{tcolorbox}[breakable,enhanced,title=\bf Instrucción en \verb"Mathematica", top=2pt,bottom=2pt,boxsep=0pt,before skip=2pt,after skip=0pt]
\begin{Verbatim}[formatcom=\letraverbatim]
{x,y}={3+(1-Sin[t])*Cos[t],(1-Sin[t])*Sin[t]};
ParametricPlot[{x,y},{t,0,2*Pi}]
ParametricPlot3D[{x*Cos[u],x*Sin[u],y},{t,0,2*Pi},
{u,0,5*Pi/3},Ticks->{{-3,-1,1,3},{-2,0,2},{-1,0,1}},
MeshStyle->Gray]
\end{Verbatim}
\end{tcolorbox}



Utilizando la función ${\bf r}(t,u)$  se hallan los siguientes vectores:

{\small
${\bf r}_{t} = \big \langle\cos(u)(\cos(2t)-\sin(t)),\sin(u)(\cos(2t)-\sin(t)),\cos(t)(1-2\sin(t)) \big\rangle.$}

${\bf r}_{u} = \big \langle\sin(u)(\cos(t)\sin(t)-3-\cos(t)),\cos(u)\cos(t)(1-\sin(t)),0\big\rangle$

y

${\bf r}_{u} = \big \langle\sin(u)(\cos(t)\sin(t)-3-\cos(t)),\cos(u)\cos(t)(1-\sin(t)),0\big\rangle.$


\begin{multline*}
{\bf r}_{t}\times {\bf r}_{u}=\bigg\langle\cos^{2}(t)(1-2\sin(t))(\sin(t)-1)\cos(u),\\
\cos(t)(1-2\sin(t))\sin (u)(\cos(t)\sin(t)-3-\cos(t)),\\
(\cos(2t)+\sin(t))(\cos(t)(\sin(t)-1)-3\sin^{2}(u))\bigg\rangle,
\end{multline*}
y con estos se obtiene el elemento diferencial de área
\begin{multline*}
\left\|{\bf r}_{t}\times {\bf r}_{u}\right\|=\frac{1}{2\sqrt{2}}\bigg|-14\cos\left(\frac{t}{2}\right)-3\cos\left(\frac{3t}{2}\right)+\cos\left(\frac{5t}{2}\right)+10\sin\left(\frac{t}{2}\right)\\
+3\sin\left(\frac{3t}{2}\right)+\sin\left(\frac{5t}{2}\right)\bigg|,
\end{multline*}
que se integra y nos proporciona el área superficial de la región en mención.
\begin{multline*}
\bigintss_{0}^{2\pi}\bigintss_{0}^{2\pi}\frac{1}{2\sqrt{2}}\bigg|-14\cos\left(\frac{t}{2}\right)-3\cos\left(\frac{3t}{2}\right)+\cos\left(\frac{5t}{2}\right)+10\sin\left(\frac{t}{2}\right)\\
+3\sin\left(\frac{3t}{2}\right)+\sin\left(\frac{5t}{2}\right)\bigg|dudt=48\pi.
\end{multline*}


Este procedimiento (y la figura de la curva revolucionada) puede evaluarse en \verb"Mathematica"
mediante la ejecución de las siguientes instrucciones:



\begin{tcolorbox}[breakable,enhanced,title=\bf Instrucción en \verb"Mathematica", top=2pt,bottom=2pt,boxsep=0pt,before skip=2pt,after skip=0pt]
\begin{Verbatim}[formatcom=\letraverbatim]
{x,y}={3+(1-Sin[t])*Cos[t],(1-Sin[t])*Sin[t]};
RevolutionPlot3D[{x,0,y},{t,0,2*Pi},MeshStyle->Gray,
RevolutionAxis->{0,0,1},PlotStyle->Orange]
r={x*Cos[u],x*Sin[u],y};
R=Norm[Cross[D[r,t],D[r,u]]];
A=Assuming[t>=0&&u>=0,Simplify[R]]
Integrate[A,{t,0,2*Pi},{u,0,2*Pi}]
\end{Verbatim}
\end{tcolorbox}




\subsubsection*{Toro cardioidico generado por \boldmath$r=1-\cos(t)$ \unboldmath}

El área de la figura de la ecuación $r=1-\cos(t)$ está dada por


{\small
\[\int_{0}^{2\pi }\int_{0}^{1-\cos \theta }rdrd\theta = \frac{3}{2}\pi\]
}
y su longitud de curva es
{\small
\[\int_{0}^{2\pi }\sqrt{\left( 1-\cos (t) \right) ^{2}+\left( \sin (t)\right)^{2}}dt=\int_{0}^{2\pi }\sqrt{2-2\cos (t) }dt= 8.\]
}


De acuerdo con la expresión (\ref{toros1}), el área superficial está determinada por
{\small
\[2\pi \int_{0}^{2\pi }\left( 3+\left( 1-\cos (t) \right) \cos
(t) \right) \sqrt{2-2\cos (t) }dt= \frac{176}{5}\pi.\]
}

El centro de masa de esta curva está dado por

{\small
\begin{eqnarray*}
\overline{x} &=& \frac{1}{8}\int_{0}^{2\pi }\left( 3+\left( 1-\cos (t) \right) \cos \left( t\right) \right) \sqrt{2-2\cos \left( t\right)}dt= \frac{11}{5}\\
\overline{y} &=& \frac{1}{8}\int_{0}^{2\pi }\left( 1-\cos \left( t\right)
\right) \sin \left( t\right) \sqrt{2-2\cos \left( t\right) }dt= 0.
\end{eqnarray*}
}


De acuerdo con el teorema de Pappus, el área de la superficie de revolución es
$2\pi \left( \frac{11}{5}\right) \left( 8\right) = \frac{176}{5}\pi .$
En este caso se usará la parametrización de la superficie:


{\small
\begin{multline*}
{\bf r}\left( t,u\right) =\big\langle( 3+\left( 1-\cos (t) \right)
\cos (t)) \cos (u) ,\left( 1-\cos (t)
\right) \cos (t) \sin \left( u\right) ,\\
\left( 1-\cos (t) \right) \sin (t) \big\rangle, \, t\in [0,2\pi],\, u\in [0,2\pi].
\end{multline*}
}

La figura obtenida es el siguiente:



\begin{figure}[H]
\centering
\caption{Toro cardiogénico construido a partir de $r=1-\cos(t)$}
\includegraphics[width=0.4\textwidth]{Img/59a.pdf} \hspace{5mm}
\includegraphics[width=0.4\textwidth]{Img/59b.pdf}

\includegraphics[width=0.4\textwidth]{Img/59c.pdf}
\fuentepropia
\end{figure}

Las instrucciones para conseguirlo se presentan a continuación.



\begin{tcolorbox}[breakable,enhanced,title=\bf Instrucción en \verb"Mathematica", top=2pt,bottom=2pt,boxsep=0pt,before skip=2pt,after skip=0pt]
\begin{Verbatim}[formatcom=\letraverbatim]
{x,y}={3+(1-Cos[t])*Cos[t],(1-Cos[t])*Sin[t]};
ParametricPlot[{x,y},{t,0,2*Pi}]
ParametricPlot3D[{x*Cos[u],x*Sin[u],y},{t,0,2*Pi},
{u,0,2*Pi}]
\end{Verbatim}
\end{tcolorbox}



Derivando parcialmente la función ${\bf r}$, hallamos los siguientes vectores:


\begin{flushleft}
\resizebox{1.0\textwidth}{!}{%
$\mathbf{r}_{t} = \left\langle (\sin(2t) - \sin t)\cos u,\; (\sin(2t) - \sin t)\sin u,\; (1 - \cos t)(2\cos t + 1) \right\rangle,$
}
\end{flushleft}
\[
{\mathbf r}_{u} = \left\langle (\cos t - 1)\cos t \sin u,\; (1 - \cos t)\cos t \cos u,\; 0 \right\rangle.
\]


Con ellos se halla el elemento diferencial de área:
\[
\left\| {\\bf r}_{t} \times {\\bf r}_{u} \right\| = \frac{1}{2} \left| -8\sin\left( \frac{t}{2} \right) - 3\sin\left( \frac{3t}{2} \right) + \sin\left( \frac{5t}{2} \right) \right|.
\]


La integración nos proporciona el área de la superficie en mención:
\[
\int_{0}^{2\pi} \int_{0}^{2\pi} \frac{1}{2} \left| -8\sin\left( \frac{t}{2} \right) - 3\sin\left( \frac{3t}{2} \right) + \sin\left( \frac{5t}{2} \right) \right| \, du \, dt = \frac{176}{5}\pi.
\]

\subsubsection*{Toro cardioidico generado por \boldmath$r=1+\cos(t)$ \unboldmath}
Al igual que en los casos anteriores, para la cardioide con ecuación
$r=1+\cos (t) ,$ el área está dada por
\[\int_{0}^{2\pi }\int_{0}^{1+\cos t}rdrdt= \frac{3}{2}\pi\]
y la longitud del arco es
\[\int_{0}^{2\pi }\sqrt{\left( 1+\cos t\right) ^{2}+\left( -\sin t\right) ^{2}}dt= 8\]
Si esta curva gira alrededor de la recta $\theta= \pi/2$, usando la expresión (\ref{toros1}),
el área superficial está dada por  
\[2\pi \int_{0}^{2\pi }\left( 3+\left( 1+\cos \left( t\right) \right) \cos
\left( t\right) \right) \sqrt{2+2\cos \left( t\right) }dt= \frac{304}{5}\pi.\]

El centro de masa de la curva está dado por
\begin{eqnarray*}
\overline{x} &=& \frac{1}{8}\int_{0}^{2\pi }\left( 3+\left( 1+\cos \left(
t\right) \right) \cos \left( t\right) \right) \sqrt{2+2\cos \left( t\right) }%
dt= \frac{19}{5} \\
\overline{y} &=& \frac{1}{8}\int_{0}^{2\pi }\left( 1+\cos \left( t\right)
\right) \sin \left( t\right) \sqrt{2+2\cos \left( t\right) }dt= 0
\end{eqnarray*}

El área de la superficie de revolución, de acuerdo con el teorema de Pappus, es  $2\pi \left( \frac{19}{5}\right) \left( 8\right) = \frac{304}{5}\pi. $


Para la superficie que se obtiene al hacer girar la figura de la curva $r=1+\cos(t)$ alrededor de la recta $\theta= \pi/2$, usaremos la parametrización.
\begin{multline*}
{\bf r}\left( t,u\right) =\bigg\langle(3+(1+\cos(t))\cos(t))\cos(u),(1+\cos(t))\cos(t)\sin(u),\\
(1+\cos(t))\sin(t)\bigg\rangle, \, t\in [0,2\pi],\, u\in [0,2\pi].
\end{multline*}

La figura obtenida es la siguiente:




\begin{figure}[H]
\centering
\caption{Toro cardiogénico construido a partir de $r=1+\cos(t)$}
\includegraphics[width=0.4\textwidth]{Img/64a.pdf} \hspace{5mm}
\includegraphics[width=0.4\textwidth]{Img/64b.pdf}

\includegraphics[width=0.4\textwidth]{Img/64c.pdf}
\fuentepropia
\end{figure}

Las siguientes instrucciones permiten obtener las figuras presentadas:



\begin{tcolorbox}[breakable,enhanced,title=\bf Instrucción en \verb"Mathematica", top=2pt,bottom=2pt,boxsep=0pt,before skip=2pt,after skip=0pt]
\begin{Verbatim}[formatcom=\letraverbatim]
x=3+{1+Cos[t]*Cos[t];
y=(1+Cos[t])*Sin[t];
ParametricPlot[{x,y},{t,0,2*Pi}]
ParametricPlot3D[{x*Cos[u],x*Sin[u],y},{t,0,2*Pi},
{u,0,4*Pi/3},MeshStyle->Gray]
\end{Verbatim}
\end{tcolorbox}



Con ayuda de la función ${\bf r}$, hallamos los siguientes vectores \linebreak
\[
\begin{split}
	\mathbf{r}_{t} &= \left\langle -\sin(t)\cos(u)(2\cos(t)+1), -\sin(t)\sin(u)(2\cos(t)+1), \right. \\
	&\quad \left. (1+\cos(t))(2\cos(t)-1) \right\rangle
\end{split}
\]
y 
\[{\bf r}_{u}= \big\langle-(3+(1+\cos (t))\cos (t))\sin (u),\linebreak (1+\cos(t))\cos(t)\cos(u),0 \big\rangle,\]
con ellos se obtiene el elemento diferencial de área
\[\left\| {\bf r}_{t}\times {\bf r}_{u}\right\| =\frac{1}{2}\left| 16\cos \left( \frac{t}{2}\right) +3\cos \left( \frac{3t}{2}\right) +\cos \left( \frac{5t}{2}\right) \right| \]
que, al integrarse, proporciona el área superficial de la región en mención.
\[\bigintss_{0}^{2\pi}\bigintss_{0}^{2\pi }\frac{1}{2}\left| 16\cos \left( \frac{t}{2}\right) +3\cos \left( \frac{3t}{2}\right) +\cos \left( \frac{5t}{2}\right) \right| dudt= \frac{304}{5}\pi.\]
Los cálculos anteriores se pueden evaluar en \verb"Mathematica" mediante la ejecución de las siguientes instrucciones:



\begin{tcolorbox}[breakable,enhanced,title=\bf Instrucción en \verb"Mathematica", top=2pt,bottom=2pt,boxsep=0pt,before skip=2pt,after skip=0pt]
\begin{Verbatim}[formatcom=\letraverbatim]
x=3+(1+Cos[t])*Cos[t];
y=(1+Cos[t])*Sin[t];
ParametricPlot[{x,0,y},{t,0,2*Pi}];RevolutionAxis->
{0,0,1}
MeshStyle->Gray,
Ticks->{{-3,0,3},{-3,0,3},
{-1,0}},
PlotStyle->Orange]
r={x*Cos[u],x*Sin[u],y};
A=Assuming[t>=0&&u>=0,
Simplify[Norm[Cross[D[r,t],D[r,u]]]]]
Integrate[A,{t,0,2*Pi},{u,0,2*Pi}]
\end{Verbatim}
\end{tcolorbox}




\begin{ejer}
Determine una expresión que defina la superficie indicada y calcule su área.
Además, determine una función que represente la curva de intersección de la superficie y halle su longitud.


\begin{enumerate}
\item La superficie sobre el paraboloide $2z=x^2+y^2$ acotada por el cilindro cardiogénico $r=1+\cos(t).$
%



\begin{figure}[H]
\centering
\caption{
	Paraboloide acotada por el cilindro cardiogénico
}
\includegraphics[width=0.4\textwidth]{Img/T_30a.pdf} \hspace{5mm}
\includegraphics[width=0.4\textwidth]{Img/T_30b.pdf}

\includegraphics[width=0.4\textwidth]{Img/T_30c.pdf}
\end{figure}

\begin{comment}
estos gráficos se obtienen ejecutando las siguientes instrucciones

\begin{Verbatim}
a=ParametricPlot3D[{(1+Cos[t])Cos[t],(1+Cos[t])*Sin[t],2*u},{t,0,2Pi},{u,0,2},
Mesh->False,PlotStyle->Orange,
PlotPoints->60];
c=ContourPlot3D[2*z==x^2+y^2,{x,-2,2},{y,-2,2},{z,0,2},
Mesh->False,ContourStyle->Cyan];
x=(1+Cos[t])Cos[t];
y=(1+Cos[t])*Sin[t];
z=(x^2+y^2)/2;
b=ParametricPlot3D[{x,y,z},{t,0,2Pi},
PlotStyle->{Black,Thickness->0.015}];
Show[c,a,b,PlotRange->{{-2,2.05},{-2,2.05},{-0.15,2}}]
\end{Verbatim}

\end{comment}

\item La superficie sobre el paraboloide hiperbólico $2z=y^2-x^2$ acotada por el cilindro cardiogénico $r=1+\cos(t).$



\begin{figure}[H]
\centering
\caption{Un paraboloide hiperbólico y un cilindro cardioidico}
\includegraphics[width=0.4\textwidth]{Img/T_31a.pdf} \hspace{5mm}
\includegraphics[width=0.4\textwidth]{Img/T_31b.pdf}
\end{figure}

\begin{comment}
estos gráficos se obtienen ejecutando las siguientes instrucciones

\begin{Verbatim}
a=ParametricPlot3D[{(1+Cos[t])Cos[t],(1+Cos[t])*Sin[t],
2*u},{t,0,2Pi},{u,-1,1},Mesh->False,
PlotStyle->Orange,PlotPoints->60];
c=ContourPlot3D[
2*z==y^2-x^2,{x,-3,3},{y,-3,3},{z,-2.5,2},
Mesh->False,ContourStyle->Blend[{Red,Green,Blue},3/4]];
x=(1+Cos[t])Cos[t];
y=(1+Cos[t])*Sin[t];
z=(y^2-x^2)/2;
b=ParametricPlot3D[{x,y,z},{t,0,2Pi},
PlotStyle->{Black,Thickness->0.015}];
Show[b,c,a,PlotRange->{{-2.5,2.5},{-2.5,2.5},{-2.1,2.6}}]
\end{Verbatim}

\end{comment}

\item La superficie sobre la esfera con ecuación $x^2+y^2+z^2=4$ acotada por el cilindro cardiogénico $1+\cos(t)$. 


\begin{figure}[H]
\centering
\caption{Esfera y cilindro cardiológico}
\includegraphics[width=0.4\textwidth]{Img/T_32a.pdf} \hspace{5mm}
\includegraphics[width=0.4\textwidth]{Img/T_32b.pdf}

\includegraphics[width=0.4\textwidth]{Img/T_32c.pdf}
\end{figure}

\begin{comment}
estos gráficos se obtienen ejecutando las siguientes instrucciones

\begin{Verbatim}
a=ParametricPlot3D[{(1+Cos[t])Cos[t],(1+Cos[t])*Sin[t],u},{t,0,2Pi},
{u,-2,2},Mesh->False,PlotStyle->Orange,PlotPoints->80];
c=ContourPlot3D[
x^2+y^2+z^2==4,{x,-2.5,2.5},{y,-2.5,2.5},{z,-2.5,2},
Mesh->False,
ContourStyle->{Blend[{Orange,Green},1/4],Opacity[0.6]},
PlotPoints->80];
x=(1+Cos[t])Cos[t];
y=(1+Cos[t])*Sin[t];
z=Sqrt[4-y^2-x^2];
b=ParametricPlot3D[{{x,y,z},{x,y,-z}},{t,0,2Pi},
PlotStyle->{Directive[Black,Thickness->0.01],
Directive[Black,Thickness->0.01]},PlotPoints->90];
Show[b,c,a,PlotRange->{{-2,2},{-2,2},{-2,2}}]
\end{Verbatim}

\end{comment}
%\clearpage

\item Un pétalo de la flor $r=\sin(3\theta)$ gira alrededor de una recta perpendicular al eje polar, que pasa por el polo.
Hallar el área de la superficie obtenida.


\begin{figure}[H]
\centering
\caption{Superficie obtenida de girar $r=\sin(3\theta)$.}
\includegraphics[width=0.4\textwidth]{Img/99a.pdf} \hspace{5mm}
\includegraphics[width=0.4\textwidth]{Img/99b.pdf}

\includegraphics[width=0.4\textwidth]{Img/99c.pdf}
\end{figure}

\begin{comment}

\begin{Verbatim}
x=Sin[3t]*Cos[t]*Cos[u];
y=Sin[3t]*Cos[t]*Sin[u];
z=Sin[3t]*Sin[t];
ParametricPlot3D[{x,z,y},{t,0,Pi/3},{u,0,13Pi/8},
Boxed->False,Axes->False,PlotTheme->"ThickSurface",
PlotStyle->Yellow,MeshStyle->Black,PlotPoints->60]
\end{Verbatim}

\end{comment}

\end{enumerate}
\end{ejer}